
\subsection{空间归一化与 Volterra 方程}
固定 $t\in I$. 在 $\Sigma\mathbin{\backslash}\{0,\pm\ii q_0\}$ 上,
以空间端态定义 Jost 解 $\chi_\pm$:
\begin{equation}\label{eq:jost-normalization}
\chi_{\pm,x}=U\chi_\pm,\qquad
\chi_\pm(x,t,z)e^{\ii\lambda(z)x\sg}\longrightarrow Y(z),
\quad x\to\pm\infty.
\end{equation}
定义去除空间振荡后的矩阵
\begin{equation}\label{eq:modified-jost}
\mu_\pm=\chi_\pm e^{\ii\lambda x\sg},\qquad
\chi_\pm=\mu_\pm e^{-\ii\lambda x\sg},\qquad
\mu_\pm\longrightarrow Y.
\end{equation}
此定义包含两个共振原像, 不需要在共振处对完整传播指数求值.
记 $\Delta Q=Q-Q_\infty$. 下面先构造分支点以外的解析列.

\paragraph{空间指数的两列与求导.}
右乘 $e^{\ii\lambda x\sg}$ 分别给第一、二列乘 $e^{\ii\lambda x}$ 与 $e^{-\ii\lambda x}$, 因而
\[
\chi_{\pm,1}=e^{-\ii\lambda x}\mu_{\pm,1},\qquad
\chi_{\pm,2}=e^{\ii\lambda x}\mu_{\pm,2}.
\]
把 $\mu_\pm=\chi_\pm e^{\ii\lambda x\sg}$ 求导,
\[
(\mu_\pm)_x
=U\chi_\pm e^{\ii\lambda x\sg}
+\chi_\pm\ii\lambda\sg e^{\ii\lambda x\sg}
=U\mu_\pm+\ii\lambda\mu_\pm\sg.
\]
矩阵 $\sg$ 可以与对角指数交换, 但不能任意穿过 $\mu_\pm$.
这里没有时间指数; $t$ 是势函数的参数, 其演化在第 \ref{ch:scattering} 节另行研究.
\begin{theorem}[Jost 解与解析列]\label{thm:jost}
设光场满足基本条件中的空间要求. 对 $\Sigma$ 上除 $0,\pm\ii q_0$ 外的谱点, Volterra 方程
\begin{equation}\label{eq:volterra}
\mu_\pm(x,z)=Y(z)+\int_{\pm\infty}^{x}
Y(z)e^{-\ii\lambda(z)(x-y)\sg}Y(z)^{-1}
\Delta Q(y)\mu_\pm(y,z)e^{\ii\lambda(z)(x-y)\sg}\,\dd y
\end{equation}
有唯一关于 $x$ 有界的解. 列 $\mu_{-,1},\mu_{+,2}$ 在 $D_+$ 解析, 列 $\mu_{+,1},\mu_{-,2}$ 在 $D_-$ 解析, 并连续延拓到各自闭域中避开 $0,\pm\ii q_0$ 的紧集. $\chi_\pm=\mu_\pm e^{-\ii\lambda x\sg}$ 是满足 \eqref{eq:jost-normalization} 的唯一空间 Jost 解, 其中包含 $z=z_a^\pm$.
\end{theorem}
\begin{proof}
修正空间方程及其背景恒等式为
\[
(\mu_\pm)_x=U\mu_\pm+\ii\lambda\mu_\pm\sg,\qquad U_\infty Y=-\ii\lambda Y\sg.
\]

\par\medskip\noindent\textbf{把变分常数公式展开成积分因子的计算.}\quad
在当前的 $z$ 上, $Y$ 与 $x$ 无关且可逆. 把
$U=U_\infty+\Delta Q$ 代入上式, 再左乘 $Y^{-1}$, 得
\begin{align*}
(Y^{-1}\mu_\pm)_x
&=Y^{-1}U_\infty\mu_\pm
+Y^{-1}\Delta Q\mu_\pm
+\ii\lambda Y^{-1}\mu_\pm\sg\\
&=-\ii\lambda\sg Y^{-1}\mu_\pm
+\ii\lambda Y^{-1}\mu_\pm\sg
+Y^{-1}\Delta Q\mu_\pm.
\end{align*}
这里用到的背景恒等式在乘以 $Y^{-1}$ 以后为
\[
Y^{-1}U_\infty Y=-\ii\lambda\sg,
\qquad
Y^{-1}U_\infty=-\ii\lambda\sg Y^{-1}.
\]
现在左右同时乘积分因子. 对整个乘积求导:
\begin{align*}
&\partial_x\bigl(e^{\ii\lambda x\sg}Y^{-1}\mu_\pm
 e^{-\ii\lambda x\sg}\bigr)\\
&=\ii\lambda\sg e^{\ii\lambda x\sg}Y^{-1}\mu_\pm
 e^{-\ii\lambda x\sg}
+e^{\ii\lambda x\sg}(Y^{-1}\mu_\pm)_x
 e^{-\ii\lambda x\sg}\\
&\quad-e^{\ii\lambda x\sg}Y^{-1}\mu_\pm
 \ii\lambda\sg e^{-\ii\lambda x\sg}\\
&=e^{\ii\lambda x\sg}Y^{-1}\Delta Q(x)\mu_\pm(x)
 e^{-\ii\lambda x\sg}.
\end{align*}
第一行求导产生的左侧项, 恰好与微分方程中的
$-\ii\lambda\sg Y^{-1}\mu_\pm$ 相消;
右侧的两项也恰好相消. 从相应的归一化端积分到 $x$, 得
\begin{align*}
&e^{\ii\lambda x\sg}Y^{-1}\mu_\pm(x)e^{-\ii\lambda x\sg}\\
&=I+\int_{\pm\infty}^{x}
 e^{\ii\lambda y\sg}Y^{-1}\Delta Q(y)\mu_\pm(y)
 e^{-\ii\lambda y\sg}\,\dd y.
\end{align*}
再左乘 $Ye^{-\ii\lambda x\sg}$, 右乘
$e^{\ii\lambda x\sg}$, 得
\begin{align*}
\mu_\pm(x)
&=Y+\int_{\pm\infty}^{x}
 Ye^{-\ii\lambda(x-y)\sg}Y^{-1}
 \Delta Q(y)\mu_\pm(y)e^{\ii\lambda(x-y)\sg}\,\dd y.
\end{align*}
对 $+\infty$ 端, 积分方向应原样保留:
\[
\int_{+\infty}^{x}(\cdots)\,\dd y
=-\int_x^{+\infty}(\cdots)\,\dd y.
\]
这个负号在求导时与下限求导的负号抵消.
在连续谱上积分因子有界, 可以从端点归一化直接得到上式.
向非实谱延拓时应逐列使用下文的有界核, 而不是要求整个
$e^{\ii\lambda x\sg}$ 在无穷远有界.

对 $Y^{-1}\mu_\pm$ 使用变分常数公式, 得到 \eqref{eq:volterra}. 令 $s_1=1,s_2=-1$, 第 $j$ 列的传播核为
\begin{equation}\label{eq:columnkernel}
K_j(r,z)=e^{\ii s_j\lambda r}Y e^{-\ii\lambda r\sg}Y^{-1},\qquad r=x-y.
\end{equation}

\par\medskip\noindent\textbf{四种列核的指数从何而来.}\quad
在矩阵积分式中取第 $j$ 列, 最右边的对角矩阵只给该列乘
$e^{\ii s_j\lambda(x-y)}$. 因此有
\[
\mu_{\pm,j}(x)
=Y_j+\int_{\pm\infty}^{x}
K_j(x-y,z)\Delta Q(y)\mu_{\pm,j}(y)\,\dd y.
\]
把两个指数相乘, 第一列和第二列分别给出
\begin{align*}
K_1(r,z)
&=Y\begin{pmatrix}e^{\ii\lambda r}e^{-\ii\lambda r}&0\\
0&e^{\ii\lambda r}e^{\ii\lambda r}\end{pmatrix}Y^{-1}
=Y\begin{pmatrix}1&0\\0&e^{2\ii\lambda r}\end{pmatrix}Y^{-1},\\
K_2(r,z)
&=Y\begin{pmatrix}e^{-\ii\lambda r}e^{-\ii\lambda r}&0\\
0&e^{-\ii\lambda r}e^{\ii\lambda r}\end{pmatrix}Y^{-1}
=Y\begin{pmatrix}e^{-2\ii\lambda r}&0\\0&1\end{pmatrix}Y^{-1}.
\end{align*}
计算模时要同时看积分方向与谱根的虚部:
\[
|e^{2\ii\lambda r}|=e^{-2r\operatorname{Im}\lambda},
\qquad
|e^{-2\ii\lambda r}|=e^{2r\operatorname{Im}\lambda}.
\]
四种情况逐一为
\begin{align*}
\mu_{-,1}:&\quad y\le x,\ r\ge0,\ \operatorname{Im}\lambda\ge0,
& e^{-2r\operatorname{Im}\lambda}&\le1,\\
\mu_{+,2}:&\quad y\ge x,\ r\le0,\ \operatorname{Im}\lambda\ge0,
& e^{2r\operatorname{Im}\lambda}&\le1,\\
\mu_{+,1}:&\quad y\ge x,\ r\le0,\ \operatorname{Im}\lambda\le0,
& e^{-2r\operatorname{Im}\lambda}&\le1,\\
\mu_{-,2}:&\quad y\le x,\ r\ge0,\ \operatorname{Im}\lambda\le0,
& e^{2r\operatorname{Im}\lambda}&\le1.
\end{align*}
前两列属于 $D_+$, 后两列属于 $D_-$.
例如 $\mu_{+,1}$ 不能沿 $D_+$ 用同一个正端积分式构造,
因为这时 $r\le0$ 而 $\operatorname{Im}\lambda>0$,
$e^{-2r\operatorname{Im}\lambda}$ 在远端增长.
这正是必须区分列而不能给整个 Jost 矩阵指定一个解析域的原因.

以 $e_1,e_2$ 表示标准基向量. 对于负端第一列, 令 $w=Y^{-1}\mu_{-,1}$, 则
\begin{equation}\label{eq:jost-column-volterra}
w(x,z)=e_1+\int_{-\infty}^{x}G(x-y,z)\Delta Q(y)Y(z)w(y,z)\,\dd y,
\end{equation}
其中
\begin{equation}\label{eq:jost-Gkernel}
G(r,z)=\diag(1,e^{2\ii\lambda r})Y^{-1}
=\dfrac{1}{\gamma(z)}
\begin{pmatrix}
1&\dfrac{\ii q_0}{z}\\[4pt]
\dfrac{\ii q_0}{z}e^{2\ii\lambda r}&e^{2\ii\lambda r}
\end{pmatrix}.
\end{equation}
第一列对应的对角指数为 $\diag(1,e^{2\ii\lambda r})$, 第二列为 $\diag(e^{-2\ii\lambda r},1)$. 由 $r=x-y$, 四列的积分方向及指数模分别满足
\begin{align*}
\mu_{-,1}:&\quad y\le x,\ r\ge0,\ \operatorname{Im}\lambda\ge0,
& e^{-2r\operatorname{Im}\lambda}&\le1,\\
\mu_{+,2}:&\quad y\ge x,\ r\le0,\ \operatorname{Im}\lambda\ge0,
& e^{2r\operatorname{Im}\lambda}&\le1,\\
\mu_{+,1}:&\quad y\ge x,\ r\le0,\ \operatorname{Im}\lambda\le0,
& e^{-2r\operatorname{Im}\lambda}&\le1,\\
\mu_{-,2}:&\quad y\le x,\ r\ge0,\ \operatorname{Im}\lambda\le0,
& e^{2r\operatorname{Im}\lambda}&\le1.
\end{align*}
故前两列对应 $D_+$, 后两列对应 $D_-$.


\par\medskip\noindent\textbf{负端第一列的两个分量积分式.}\quad
$w=Y^{-1}\mu_{-,1}$ 沿用原文记号. 先把进入积分的矩阵乘法写全:
\[
\Delta Q(y)=\begin{pmatrix}0&q(y)-q_0\\
-\conj{q(y)}+q_0&0\end{pmatrix},
\qquad
Yw=\begin{pmatrix}
w_1-\dfrac{\ii q_0}{z}w_2\\
-\dfrac{\ii q_0}{z}w_1+w_2
\end{pmatrix}.
\]
因此
\[
\Delta Q(y)Yw
=\begin{pmatrix}
(q-q_0)\left(-\dfrac{\ii q_0}{z}w_1+w_2\right)\\
-(\conj{q}-q_0)\left(w_1-\dfrac{\ii q_0}{z}w_2\right)
\end{pmatrix}.
\]
在这里所有 $q,w$ 均取 $y$ 处的值. 再左乘 $G$, 得
\begin{align*}
w_1(x)
&=1+\dfrac{1}{\gamma}\int_{-\infty}^{x}
(q-q_0)\left(-\dfrac{\ii q_0}{z}w_1+w_2\right)\,\dd y\\
&\quad-\dfrac{\ii q_0}{z\gamma}\int_{-\infty}^{x}(\conj{q}-q_0)
\left(w_1-\dfrac{\ii q_0}{z}w_2\right)\,\dd y,\\
w_2(x)
&=\dfrac{\ii q_0}{z\gamma}\int_{-\infty}^{x}e^{2\ii\lambda(x-y)}(q-q_0)
\left(-\dfrac{\ii q_0}{z}w_1+w_2\right)\,\dd y\\
&\quad-\dfrac{1}{\gamma}\int_{-\infty}^{x}e^{2\ii\lambda(x-y)}
(\conj{q}-q_0)\left(w_1-\dfrac{\ii q_0}{z}w_2\right)\,\dd y.
\end{align*}
这两个式子显示: 背景列 $e_1$ 的第二分量为零,
而唯一可能增长的指数只出现在第二分量的传播中.
同样的矩阵乘法对第二列将指数置于第一行.

采用向量范数 $\|v\|_1=|v_1|+|v_2|$ 及其诱导矩阵范数. 由
\[
\|\Delta Q(y)\|_1=|q(y)-q_0|,\qquad
\|Y\|_1=1+\dfrac{q_0}{|z|},\qquad
\|Y^{-1}\|_1=\dfrac{1+\dfrac{q_0}{|z|}}{|\gamma(z)|},
\]
可得 HMB 积分核的具体估计
\begin{equation}\label{eq:jost-kernel-estimate}
\|G(x-y,z)\Delta Q(y)Y\|_1
\le c(z)|q(y)-q_0|,\qquad
c(z)=\dfrac{\left(1+\dfrac{q_0}{|z|}\right)^2}{\left|1+\dfrac{q_0^2}{z^2}\right|}.
\end{equation}
取 $K\subset\overline{D_+}\mathbin{\backslash}\{0,\pm\ii q_0\}$ 为紧集, 记
\[
c_{K}=\max\limits_{z\in K}c(z),\qquad
M_{K}(x)=c_{K}\int_{-\infty}^{x}|q(y)-q_0|\,\dd y.
\]

\par\medskip\noindent\textbf{矩阵范数估计中的每一个因子.}\quad
诱导 $1$ 范数是各列绝对值之和的最大值. 由上面三个矩阵的列和,
\begin{align*}
\|\Delta Q\|_1
&=\max\{|\conj{q}-q_0|,|q-q_0|\}=|q-q_0|,\\
\|Y\|_1
&=\max\left\{1+\dfrac{q_0}{|z|},1+\dfrac{q_0}{|z|}\right\}
=1+\dfrac{q_0}{|z|},\\
\|G(r,z)\|_1
&\le\dfrac{1}{|\gamma|}
\max\left\{1+\dfrac{q_0}{|z|},\dfrac{q_0}{|z|}+1\right\}
=\dfrac{1+\dfrac{q_0}{|z|}}{|\gamma|}.
\end{align*}
最后一行使用了允许方向上的 $|e^{2\ii\lambda r}|\le1$.
接着按乘法次序估计,
\begin{align*}
\|G\Delta QYw\|_1
&\le\|G\|_1\|\Delta Q\|_1\|Y\|_1\|w\|_1\\
&\le\dfrac{\left(1+\dfrac{q_0}{|z|}\right)^2}{|\gamma|}
|q-q_0|\|w\|_1.
\end{align*}
紧集避开 $z=0$ 与 $\gamma=0$ 的两点, 因而 $c_{ K}$ 有限.
$M_{ K}$ 是非减的绝对连续函数, 且几乎处处满足
\[
M_{ K}'(x)=c_{ K}|q(x)-q_0|,
\qquad M_{ K}(-\infty)=0.
\]
这一个导数等式正是下文阶乘分母的来源.

构造 Neumann 级数
\begin{equation}\label{eq:jost-neumann}
\begin{aligned}
w&=\sum\limits_{n=0}^{\infty}w^{(n)},\qquad w^{(0)}=e_1,\\
w^{(n+1)}&=\int_{-\infty}^{x}G(x-y,z)\Delta Q(y)Yw^{(n)}(y,z)\,\dd y.
\end{aligned}
\end{equation}
若 $\|w^{(n)}(y,z)\|_1\le \dfrac{M_{K}(y)^n}{n!}$, 则
\[
\|w^{(n+1)}(x,z)\|_1\le\dfrac{1}{n!}\int_{-\infty}^{x}c_{K}|q(y)-q_0|M_{K}(y)^n\,\dd y
=\dfrac{M_{K}(x)^{n+1}}{(n+1)!}.
\]
归纳即得
\begin{equation}\label{eq:jost-factorial}
\|w^{(n)}(x,z)\|_1\le\dfrac{M_{K}(x)^n}{n!},\qquad
\|w(x,z)-e_1\|_1\le e^{M_{K}(x)}-1.
\end{equation}

\par\medskip\noindent\textbf{把逐次积分真正展开到阶乘.}\quad
第一次与第二次迭代分别满足
\begin{align*}
\|w^{(1)}(x,z)\|_1
&\le\int_{-\infty}^{x}M_{ K}'(y)\,\dd y
=M_{ K}(x),\\
\|w^{(2)}(x,z)\|_1
&\le\int_{-\infty}^{x}M_{ K}'(y)
\int_{-\infty}^{y}M_{ K}'(y_1)\,\dd y_1\,\dd y\\
&=\int_{-\infty}^{x}M_{ K}'(y)M_{ K}(y)\,\dd y
=\dfrac{M_{ K}(x)^2}{2}.
\end{align*}
第三次同样得到
\[
\|w^{(3)}(x,z)\|_1
\le\int_{-\infty}^{x}M_{ K}'(y)
\dfrac{M_{ K}(y)^2}{2}\,\dd y
=\dfrac{M_{ K}(x)^3}{6}.
\]
一般一步的等号由微积分基本定理逐行给出:
\begin{align*}
\dfrac{1}{n!}\int_{-\infty}^{x}
M_{ K}'(y)M_{ K}(y)^n\,\dd y
&=\dfrac{1}{n!}\int_{-\infty}^{x}
\dfrac{1}{n+1}\partial_y\bigl(M_{ K}(y)^{n+1}\bigr)\,\dd y\\
&=\dfrac{M_{ K}(x)^{n+1}-M_{ K}(-\infty)^{n+1}}{(n+1)!}\\
&=\dfrac{M_{ K}(x)^{n+1}}{(n+1)!}.
\end{align*}
等价地, 第 $n$ 次迭代的积分区域是
\[
-\infty<y_n<y_{n-1}<\cdots<y_1<x.
\]
把 $( -\infty,x)^n$ 按 $n$ 个坐标的先后次序分成 $n!$ 个区域;
被积的范数上界是各坐标中同一非负函数的乘积, 因而每个区域的积分相等.
于是其积分为整个乘积积分的 $\dfrac{1}{n!}$, 即
$\dfrac{M_{ K}(x)^n}{n!}$.
求和时从 $n=1$ 开始的部分为
\[
\sum\limits_{n=1}^{\infty}\dfrac{M_{ K}(x)^n}{n!}
=e^{M_{ K}(x)}-1.
\]
特别地, 不要求 $\|q-q_0\|_{L^1}$ 小于某个数;
阶乘使任意有限的 $L^1$ 质量都能给出收敛.

由于 $M_{K}(x)\le c_{K}\|q-q_0\|_{L^1}$, 级数对实 $x$ 及 $z\in K$ 绝对一致收敛, 且在 $x\to-\infty$ 时趋于 $e_1$. 每一迭代项在 $D_+$ 解析, 并在相应闭域上连续; 上述一致估计保证其和 $w$ 及 $\mu_{-,1}=Yw$ 也具有这些性质.


\par\medskip\noindent\textbf{参数解析性与闭域连续性所用的控制.}\quad
对任一固定有限的 $x$, 先把积分在 $y=-R$ 处截断.
在有限区间上, 核的各矩阵元都是 $z$ 的解析函数,
因为 $Y,Y^{-1},\lambda$ 在当前紧集附近解析, 指数是整函数.
再让 $R\to+\infty$, 积分尾部的范数至多为
\[
c_{ K}\sup\limits_{y,z}\|w^{(n)}(y,z)\|_1
\int_{-\infty}^{-R}|q(y)-q_0|\,\dd y,
\]
它与 $z\in K$ 无关并趋于零.
所以每一项 $w^{(n)}$ 都是解析函数的一致极限.
对所有迭代项再使用
\[
\sum\limits_{n=0}^{\infty}
\sup\limits_{x\in\mathbb R,z\in K}\|w^{(n)}(x,z)\|_1
\le\sum\limits_{n=0}^{\infty}
\dfrac{\bigl(c_{ K}\|q-q_0\|_{L^1}\bigr)^n}{n!}<\infty.
\]
解析函数列的紧集一致极限仍然解析.
若 $z$ 从解析域趋向闭域边界, 核逐点连续,
上面同一个可积上界仍适用; 对每项与级数分别取极限,
即可交换边界极限、积分和无穷求和.
这里不声称仅靠 $L^1$ 就能在边界上任意次求 $z$ 导数;
对导数直接估计会出现 $|x-y|$, 需要与所求结论相适应的矩条件.

对两个有界解之差 $d$ 迭代 $n$ 次, 有
\[
\|d(x,z)\|_1\le\|d(\cdot,z)\|_{L^\infty}\dfrac{M_{K}(x)^n}{n!}\longrightarrow0,
\]
从而解唯一. 其余三列同理.

\par\medskip\noindent\textbf{唯一性与空间微分各自怎样完成.}\quad
差 $d$ 的自由项为零. 第一次应用积分方程给出
\[
\|d(x,z)\|_1\le\|d\|_{L^\infty}M_{ K}(x).
\]
把同一个积分式再次代入积分中的 $d(y,z)$,
积分区域成为 $y_1\le x,y_2\le y_1$, 故
\[
\|d(x,z)\|_1\le\|d\|_{L^\infty}
\dfrac{M_{ K}(x)^2}{2!}.
\]
重复 $n$ 次后就是原文所列估计. 对固定 $M_{ K}(x)$,
相邻上界之比为 $\dfrac{M_{ K}(x)}{n+1}$, 趋于零,
因此右端趋于零, 从而 $d(x,z)=0$.
对空间导数, 直接使用原文的列核有
\[
K_j(0,z)=I,
\qquad
\partial_rK_j(r,z)=(U_\infty+\ii s_j\lambda I)K_j(r,z).
\]
这是因为 $Ye^{-\ii\lambda r\sg}Y^{-1}=e^{rU_\infty}$.
逐列对积分式求导, 可得
\begin{align*}
(\mu_{\pm,j})_x
&=\Delta Q(x)\mu_{\pm,j}(x)
+\int_{\pm\infty}^{x}(U_\infty+\ii s_j\lambda I)
 K_j(x-y,z)\Delta Q(y)\mu_{\pm,j}(y)\,\dd y\\
&=\Delta Q\mu_{\pm,j}
+(U_\infty+\ii s_j\lambda I)(\mu_{\pm,j}-Y_j)\\
&=(U+\ii s_j\lambda I)\mu_{\pm,j}.
\end{align*}
最后使用 $(U_\infty+\ii s_j\lambda I)Y_j=0$.
故两列合起来满足修正空间方程.
再计算恢复指数后的导数:
\begin{align*}
\partial_x(\mu_{\pm,j}e^{\ii s_j(-\lambda x)})
&=(U\mu_{\pm,j}+\ii s_j\lambda\mu_{\pm,j})e^{\ii s_j(-\lambda x)}
-\ii s_j\lambda\mu_{\pm,j}e^{\ii s_j(-\lambda x)}\\
&=U\chi_{\pm,j}.
\end{align*}
在目前的低正则性假设下, 上述空间导数理解为几乎处处导数,
解在有限区间上绝对连续; 若 $q$ 连续, 等式处处成立.
\end{proof}

\subsection{分支点的连续延拓}
\begin{lemma}[分支点的空间极限]\label{prop:jost-branch}
基本条件给出加权可积性
\begin{equation}\label{eq:jost-weighted}
\int_{-\infty}^{+\infty}(1+|x|)|q(x)-q_0|\,\dd x<\infty.
\end{equation}
在此条件下, 每个 Jost 解析列均可从其闭解析域连续延拓到分支点 $z_b=\pm\ii q_0$. 端点值满足
\begin{equation}\label{eq:jost-branch-volterra}
\mu_\pm(x,z_b)=Y(z_b)+\int_{\pm\infty}^{x}
[I+(x-y)U_\infty(z_b)]\Delta Q(y)\mu_\pm(y,z_b)\,\dd y,
\end{equation}
并在对应归一化端趋于 $Y(z_b)$, 满足空间谱方程, 且秩为一.
\end{lemma}
\begin{proof}

\par\medskip\noindent\textbf{为什么分支点必须重新估计核.}\quad
当 $z\to z_b$, $\gamma(z)\to0$, 因而前面的
$c(z)$ 发散. 这只说明将 $Y$ 与 $Y^{-1}$ 分开估计过粗,
并不说明二者夹住指数后的整体发散.
先把背景矩阵的平方完整相乘:
\[
\begin{pmatrix}-\ii k&q_0\\-q_0&\ii k\end{pmatrix}^2
=\begin{pmatrix}
-k^2-q_0^2&-\ii kq_0+\ii kq_0\\
\ii kq_0-\ii kq_0&-q_0^2-k^2
\end{pmatrix}
=-\lambda^2I.
\]
所以偶数次幂与奇数次幂分别为
\[
U_\infty^{2n}=(-1)^n\lambda^{2n}I,
\qquad
U_\infty^{2n+1}=(-1)^n\lambda^{2n}U_\infty.
\]
代入矩阵指数的幂级数并把偶奇两部分分开,
\begin{align*}
e^{rU_\infty}
&=\sum\limits_{n=0}^{\infty}\dfrac{(-1)^n(\lambda r)^{2n}}{(2n)!}I
+\sum\limits_{n=0}^{\infty}\dfrac{(-1)^n\lambda^{2n}r^{2n+1}}{(2n+1)!}U_\infty\\
&=\cos(\lambda r)I+\dfrac{\sin(\lambda r)}{\lambda}U_\infty.
\end{align*}
在 $\lambda=0$ 处第二个系数按连续值取 $r$.

由 $U_\infty^2=-\lambda^2I$, 背景传播矩阵可写为
\begin{equation}\label{eq:background-propagator}
e^{rU_\infty}
=\cos(\lambda r)I+\dfrac{\sin(\lambda r)}{\lambda}U_\infty
\longrightarrow I+rU_\infty(z_b),\qquad z\to z_b.
\end{equation}
将 \eqref{eq:columnkernel} 写为
\begin{equation}\label{eq:branch-columnkernel}
K_j(r,z)=\dfrac{1+e^{2\ii s_j\lambda r}}{2}I
+\dfrac{e^{2\ii s_j\lambda r}-1}{2\ii s_j\lambda}U_\infty(z).
\end{equation}

\par\medskip\noindent\textbf{消掉分支点的表面分母.}\quad
将上式乘上列核中的 $e^{\ii s_j\lambda r}$.
利用 $s_j=\pm1$, 分别计算
\begin{align*}
e^{\ii s_j\lambda r}\cos(\lambda r)
&=\dfrac{1+e^{2\ii s_j\lambda r}}{2},\\
e^{\ii s_j\lambda r}\dfrac{\sin(\lambda r)}{\lambda}
&=\dfrac{e^{2\ii s_j\lambda r}-1}{2\ii s_j\lambda}.
\end{align*}
这得到原文的列核表达式. 第二个系数不是实质的极点:
\begin{align*}
e^{2\ii s_j\lambda r}-1
&=\int_0^1\partial_\tau
 e^{2\ii s_j\tau\lambda r}\,\dd\tau\\
&=2\ii s_j\lambda r\int_0^1
 e^{2\ii s_j\tau\lambda r}\,\dd\tau.
\end{align*}
当 $\lambda\to0$ 时, 这个积分趋于 1, 因而商趋于 $r$.
如果 $s_j\operatorname{Im}\lambda\,r\ge0$, 则每个
$0\le\tau\le1$ 上的指数模均不超过 1, 于是
\[
\left|\dfrac{1+e^{2\ii s_j\lambda r}}{2}\right|\le1,
\qquad
\left|\dfrac{e^{2\ii s_j\lambda r}-1}{2\ii s_j\lambda}\right|
\le|r|.
\]
故只需控制 $U_\infty$ 在分支点邻域的范数,
就能得到 $1+C_b|r|$ 的一致上界.
这解释了一阶空间矩为何足以处理分支点.

在该列允许的积分方向及闭解析域内, $s_j\operatorname{Im}\lambda\,r\ge0$, 且
\[
\dfrac{e^{2\ii s_j\lambda r}-1}{2\ii s_j\lambda}
=r\int_0^1 e^{2\ii s_j\tau\lambda r}\,\dd\tau,\qquad
\|K_j(r,z)\|_1\le1+C_b|r|.
\]
其中 $C_b$ 在 $z_b$ 的一个闭邻域内一致有界, 且 $K_j(r,z_b)=I+rU_\infty(z_b)$.

以负端为例, 固定 $b\in\mathbb R$, 对 $y\le x\le b$ 有
\[
1+C_b|x-y|\le A_b(1+|y|),\qquad
L_b(x)=A_b\int_{-\infty}^{x}(1+|y|)|q(y)-q_0|\,\dd y.
\]

\par\medskip\noindent\textbf{加权估计与半直线归一化.}\quad
当 $y\le x\le b$ 时,
\[
0\le x-y\le b-y\le|b|+|y|,
\]
故
\[
1+C_b|x-y|
\le 1+C_b|b|+C_b|y|
\le (1+C_b|b|+C_b)(1+|y|).
\]
因而可以把 $A_b$ 选为括号中的常数或更大的常数.
这不是对整个实轴统一控制 $|x-y|$;
它是在任意固定半直线 $x\le b$ 上控制积分核.
对第 $n$ 次迭代应用与前面完全相同的有序积分,
每个 $\Delta Q$ 的范数与 $1+C_b|x-y|$ 相乘后,
都由 $A_b(1+|y|)|q(y)-q_0|$ 控制.
于是阶乘估计的积分函数变为 $L_b$, 不再出现 $\dfrac{1}{\gamma}$.
又有
\[
\|\mu_{-,j}(x,z)-Y_j(z)\|_1
\le\left(\sup\limits_z\|Y_j(z)\|_1\right)(e^{L_b(x)}-1).
\]
加权可积性给出 $L_b(x)\to0$ 当 $x\to-\infty$,
所以端点归一化与 $z\to z_b$ 极限可以兼容.
在正端, 对 $b\le x\le y$ 用
$y-x\le |y|+|b|$, 并把积分上界改为从 $x$ 到 $+\infty$,
其余每一步保持相同.

列积分方程的第 $n$ 次迭代满足
\[
\|\mu_{-,j}^{(n)}(x,z)\|_1
\le C_Y\dfrac{L_b(x)^n}{n!},\qquad x\le b,
\]
其中 $C_Y$ 为该邻域内 $\|Y_j(z)\|_1$ 的统一界. 加权条件保证级数一致收敛, 逐项取极限得到 \eqref{eq:jost-branch-volterra}; $L_b(x)\to0$ 给出负端归一化. 正端的反向积分同理, 对两解之差的迭代也给出端点唯一性.


\par\medskip\noindent\textbf{端点空间方程与秩一结论的具体计算.}\quad
在 $z_b=\ii q_0$ 处,
\[
Y(z_b)=\begin{pmatrix}1&-1\\-1&1\end{pmatrix},
\qquad Y_2(z_b)=-Y_1(z_b).
\]
在 $z_b=-\ii q_0$ 处,
\[
Y(z_b)=\begin{pmatrix}1&1\\1&1\end{pmatrix},
\qquad Y_2(z_b)=Y_1(z_b).
\]
这两种线性关系都能逐列代入同一个端点 Volterra 方程.
负号情形下, $\mu_{\pm,2}+\mu_{\pm,1}$ 的自由项为零;
正号情形下, $\mu_{\pm,2}-\mu_{\pm,1}$ 的自由项为零.
由该半直线上的唯一性, 相应组合恒为零.
求端点方程的空间导数时,
\begin{align*}
(\mu_\pm)_x
&=\Delta Q(x)\mu_\pm(x)
+U_\infty\int_{\pm\infty}^{x}
\Delta Q(y)\mu_\pm(y)\,\dd y,\\
U_\infty\mu_\pm(x)
&=U_\infty Y
+\int_{\pm\infty}^{x}
\bigl(U_\infty+(x-y)U_\infty^2\bigr)
\Delta Q(y)\mu_\pm(y)\,\dd y\\
&=U_\infty\int_{\pm\infty}^{x}\Delta Q(y)\mu_\pm(y)\,\dd y.
\end{align*}
因此 $(\mu_\pm)_x=(U_\infty+\Delta Q)\mu_\pm=U\mu_\pm$.
第一列在归一化端趋于上面非零的 $Y_1$.
若它在某个有限点为零, 一阶线性初值问题的唯一性会使它处处为零,
与非零端态矛盾. 所以两列始终相关且至少一列非零, 矩阵秩恰为一.

对 \eqref{eq:jost-branch-volterra} 求导得到 $(\mu_\pm)_x=U(z_b)\mu_\pm$. 由端点唯一性, $Y(z_b)$ 两列的线性关系传递给 $\mu_\pm(x,z_b)$. 归一化端值非零, 线性系统的非零解不能在有限点为零, 故其秩恒为一.
\end{proof}

共振原像 $z_a^\pm$ 不使空间积分核退化, 因而 $\mu_\pm$ 在此仍由 \eqref{eq:volterra} 正则取值. 修正时间方程及共振相容性见第 \ref{ch:scattering} 节.

\subsection{散射矩阵与解析系数}
记 $\Sigma^\circ=\Sigma\mathbin{\backslash}\{0,\pm\ii q_0\}$. 在 $\Sigma^\circ$ 上, 两组 Jost 解为同一空间谱方程的基本解矩阵, 因而存在与 $x$ 无关的散射矩阵
\begin{equation}\label{eq:Sdef}
\chi_+(x,t,z)=\chi_-(x,t,z)S(z,t),\qquad S(z,t)=(s_{jk}(z,t))_{j,k=1}^{2}.
\end{equation}

\paragraph{为什么两组基底只相差一个常矩阵.}
先在非分支点处讨论. 下面将证明 $\det\chi_-=\gamma\ne0$,
所以散射关系等价于 $S=\chi_-^{-1}\chi_+$.
对恒等式 $\chi_-^{-1}\chi_-=I$ 求导,
\[
(\chi_-^{-1})_x\chi_-+\chi_-^{-1}(\chi_-)_x=0,
\qquad
(\chi_-^{-1})_x=-\chi_-^{-1}U.
\]
于是
\begin{align*}
S_x
&=(\chi_-^{-1})_x\chi_++\chi_-^{-1}(\chi_+)_x\\
&=-\chi_-^{-1}U\chi_++\chi_-^{-1}U\chi_+=0.
\end{align*}
这里 $U$ 可以随 $x$ 变化; 只因两组矩阵满足同一个空间方程,
这两个包含 $U$ 的项才精确相消.
把右乘矩阵逐元写出,
\[
(\chi_{-,1},\chi_{-,2})
\begin{pmatrix}s_{11}&s_{12}\\s_{21}&s_{22}\end{pmatrix}
=\bigl(s_{11}\chi_{-,1}+s_{21}\chi_{-,2},
 s_{12}\chi_{-,1}+s_{22}\chi_{-,2}\bigr).
\]
注意第一列的两个系数是 $s_{11},s_{21}$,
而不是散射矩阵的第一行. 所有后续反射系数的符号都依赖这一约定.

以下省略散射系数的 $t$ 自变量; 本节只证明它们与 $x$ 无关. 其列关系为
\begin{equation}\label{eq:scattering-columns}
\chi_{+,1}=s_{11}\chi_{-,1}+s_{21}\chi_{-,2},\qquad
\chi_{+,2}=s_{12}\chi_{-,1}+s_{22}\chi_{-,2}.
\end{equation}

\begin{proposition}[散射矩阵与 Wronskian]\label{prop:scattering}
在定理 \ref{thm:jost} 的条件下,
\begin{equation}\label{eq:detjost}
\det\chi_\pm=\det\mu_\pm=\gamma(z),\qquad \det S=1.
\end{equation}
散射系数满足
\begin{align}
s_{11}&=\dfrac{\det(\chi_{+,1},\chi_{-,2})}{\gamma},&
s_{22}&=\dfrac{\det(\chi_{-,1},\chi_{+,2})}{\gamma},\label{eq:wronskian}\\
s_{21}&=\dfrac{\det(\chi_{-,1},\chi_{+,1})}{\gamma},&
s_{12}&=\dfrac{\det(\chi_{+,2},\chi_{-,2})}{\gamma}.\nonumber
\end{align}
其中 $s_{11}$ 在 $D_-$ 解析, $s_{22}$ 在 $D_+$ 解析, 并连续延拓到相应闭域中避开 $0,\pm\ii q_0$ 的紧集. $s_{12},s_{21}$ 在 $\Sigma^\circ$ 上连续; 在共振原像处使用空间归一化值.
\end{proposition}
\begin{proof}

\par\medskip\noindent\textbf{Abel 恒等式在这个二阶系统中的完整求导.}\quad
空间方程的四个分量是
\begin{align*}
(\chi_{11})_x&=-\ii k\chi_{11}+q\chi_{21},&
(\chi_{12})_x&=-\ii k\chi_{12}+q\chi_{22},\\
(\chi_{21})_x&=-\conj{q}\chi_{11}+\ii k\chi_{21},&
(\chi_{22})_x&=-\conj{q}\chi_{12}+\ii k\chi_{22}.
\end{align*}
这里 $\chi$ 可以逐次取 $\chi_+$ 或 $\chi_-$,
并未引入新的谱解. 对二阶行列式的两项分别求导:
\begin{align*}
\partial_x\det\chi
&=\partial_x(\chi_{11}\chi_{22}-\chi_{12}\chi_{21})\\
&=(-\ii k\chi_{11}+q\chi_{21})\chi_{22}
+\chi_{11}(-\conj{q}\chi_{12}+\ii k\chi_{22})\\
&\quad-(-\ii k\chi_{12}+q\chi_{22})\chi_{21}
-\chi_{12}(-\conj{q}\chi_{11}+\ii k\chi_{21})\\
&=-\ii k\chi_{11}\chi_{22}+\ii k\chi_{11}\chi_{22}
+\ii k\chi_{12}\chi_{21}-\ii k\chi_{12}\chi_{21}\\
&\quad+q\chi_{21}\chi_{22}-q\chi_{22}\chi_{21}
-\conj{q}\chi_{11}\chi_{12}+\conj{q}\chi_{12}\chi_{11}=0.
\end{align*}
每个矩阵元是标量, 所以最后两组也能交换次序而相消.
归一化端的行列式为
\begin{align*}
\det\chi_\pm
&=\det\mu_\pm\det e^{-\ii\lambda x\sg}
=\det\mu_\pm\,e^{-\ii\lambda x}e^{\ii\lambda x}
=\det\mu_\pm,\\
\lim\limits_{x\to\pm\infty}\det\mu_\pm
&=\det\begin{pmatrix}1&-\dfrac{\ii q_0}{z}\\
-\dfrac{\ii q_0}{z}&1\end{pmatrix}
=1-\left(-\dfrac{\ii q_0}{z}\right)^2
=1+\dfrac{q_0^2}{z^2}=\gamma.
\end{align*}
由于行列式不随 $x$ 变化, 它等于自己的端点极限.
对散射关系取行列式,
\[
\gamma=\gamma\det S,
\qquad \gamma\ne0,
\qquad \det S=1.
\]

由 $\tr U=0$, Abel 恒等式给出 $\partial_x\det\chi_\pm=0$. 在归一化端取极限, 并使用 $\det e^{-\ii\lambda x\sg}=1$, 得到 $\det\chi_\pm=\gamma$. 对 \eqref{eq:Sdef} 取行列式即得 $\det S=1$.
\par\medskip\noindent\textbf{四个 Wronskian 分别怎样消去另一项.}\quad
二列行列式对每列都是线性的, 且两列相同的行列式为零.
对第一条列关系, 两种取法为
\begin{align*}
\det(\chi_{+,1},\chi_{-,2})
&=s_{11}\det(\chi_{-,1},\chi_{-,2})
+s_{21}\det(\chi_{-,2},\chi_{-,2})\\
&=s_{11}\gamma+s_{21}\cdot0=s_{11}\gamma,\\
\det(\chi_{-,1},\chi_{+,1})
&=s_{11}\det(\chi_{-,1},\chi_{-,1})
+s_{21}\det(\chi_{-,1},\chi_{-,2})\\
&=s_{11}\cdot0+s_{21}\gamma=s_{21}\gamma.
\end{align*}
对第二条列关系, 应按下面的列次序取行列式:
\begin{align*}
\det(\chi_{-,1},\chi_{+,2})
&=s_{12}\det(\chi_{-,1},\chi_{-,1})
+s_{22}\det(\chi_{-,1},\chi_{-,2})\\
&=s_{22}\gamma,\\
\det(\chi_{+,2},\chi_{-,2})
&=s_{12}\det(\chi_{-,1},\chi_{-,2})
+s_{22}\det(\chi_{-,2},\chi_{-,2})\\
&=s_{12}\gamma.
\end{align*}
例如, 若把最后一个行列式的两列反过来,
会得到 $-s_{12}\gamma$, 因而不能随意交换列序.
除以 $\gamma$ 后恰好得到原文四个散射系数.

对第一条列关系分别与 $\chi_{-,2}$、$\chi_{-,1}$ 取行列式, 有
\[
\det(\chi_{+,1},\chi_{-,2})=s_{11}\gamma,\qquad
\det(\chi_{-,1},\chi_{+,1})=s_{21}\gamma.
\]
第二条列关系同理给出 $s_{12},s_{22}$, 即 \eqref{eq:wronskian}.
\par\medskip\noindent\textbf{把四个行列式中的指数也逐个抽出.}\quad
由第一列与第二列分别带 $e^{-\ii\lambda x}$ 与 $e^{\ii\lambda x}$,
\begin{align*}
\det(\chi_{+,1},\chi_{-,2})
&=e^{-\ii\lambda x}e^{\ii\lambda x}
\det(\mu_{+,1},\mu_{-,2})
=\det(\mu_{+,1},\mu_{-,2}),\\
\det(\chi_{-,1},\chi_{+,2})
&=e^{-\ii\lambda x}e^{\ii\lambda x}
\det(\mu_{-,1},\mu_{+,2})
=\det(\mu_{-,1},\mu_{+,2}),\\
\det(\chi_{-,1},\chi_{+,1})
&=e^{-2\ii\lambda x}\det(\mu_{-,1},\mu_{+,1}),\\
\det(\chi_{+,2},\chi_{-,2})
&=e^{2\ii\lambda x}\det(\mu_{+,2},\mu_{-,2}).
\end{align*}
因此四个系数也可写为
\begin{align*}
s_{11}&=\dfrac{(\mu_+)_{11}(\mu_-)_{22}
-(\mu_-)_{12}(\mu_+)_{21}}{\gamma},\\
s_{22}&=\dfrac{(\mu_-)_{11}(\mu_+)_{22}
-(\mu_+)_{12}(\mu_-)_{21}}{\gamma},\\
s_{21}&=\dfrac{e^{-2\ii\lambda x}}{\gamma}
\bigl((\mu_-)_{11}(\mu_+)_{21}
-(\mu_+)_{11}(\mu_-)_{21}\bigr),\\
s_{12}&=\dfrac{e^{2\ii\lambda x}}{\gamma}
\bigl((\mu_+)_{12}(\mu_-)_{22}
-(\mu_-)_{12}(\mu_+)_{22}\bigr).
\end{align*}
$s_{22}$ 的两列恰为 $\mu_{-,1},\mu_{+,2}$,
二者都在 $D_+$ 解析; $s_{11}$ 的两列恰为
$\mu_{+,1},\mu_{-,2}$, 二者都在 $D_-$ 解析.
$\gamma$ 在相应开域没有零点, 所以除法保留解析性.
另一方面, $s_{21}$ 与 $s_{12}$ 所用两列来自相反解析域,
这里只能在共同边界上取值, 不能由这两个式子直接宣称全域解析.

对透射系数, 两列的指数相消, 故
\[
s_{11}=\dfrac{\det(\mu_{+,1},\mu_{-,2})}{\gamma},\qquad
s_{22}=\dfrac{\det(\mu_{-,1},\mu_{+,2})}{\gamma}.
\]
两个分子分别由 $D_-$ 和 $D_+$ 的解析列组成, 且 $\gamma$ 的零点只位于分支点. 定理 \ref{thm:jost} 因而给出所述解析性与连续性. 对非对角元, 列指数只留下连续因子 $e^{\pm2\ii(-\lambda x)}$, 故它们在 $\Sigma^\circ$ 上连续.
\end{proof}

按照 $\chi_+=\chi_-S$ 的约定, 以下称 $s_{11},s_{22}$ 为透射系数. 在其非零的谱弧上定义反射系数
\begin{equation}\label{eq:reflection}
\rho(z)=\dfrac{s_{21}(z)}{s_{11}(z)},\qquad
\widetilde\rho(z)=\dfrac{s_{12}(z)}{s_{22}(z)}.
\end{equation}
由 $\det S=1$,
\begin{equation}\label{eq:scattering-determinant}
s_{11}s_{22}-s_{12}s_{21}=1,\qquad
s_{11}s_{22}(1-\rho\widetilde\rho)=1.
\end{equation}

\paragraph{反射系数与行列式约束的代数.}
由定义, 在分母不为零的谱弧上有
\[
s_{21}=\rho s_{11},\qquad s_{12}=\widetilde\rho s_{22}.
\]
将这两式原次序代入行列式,
\begin{align*}
s_{11}s_{22}-s_{12}s_{21}
&=s_{11}s_{22}-(\widetilde\rho s_{22})(\rho s_{11})\\
&=s_{11}s_{22}(1-\rho\widetilde\rho)=1.
\end{align*}
由于这里都是标量函数, 可以交换乘法次序.
这也说明 $1-\rho\widetilde\rho$ 不能在该谱弧上为零.
共轭约化之后, 实轴上会有
$\widetilde\rho=-\conj\rho$, 从而该因子变成
$1+|\rho|^2$; 这里先保留原有符号形式,
以免把只在实轴成立的等式用于圆谱的非实点.

后续矩阵 RHP 取普通连续谱上 $s_{11}s_{22}\ne0$ 的散射数据; 分支点的边界正则性另行列入 RHP 的条件.

\subsection{共轭与换叶约化}\label{ch:symmetry}
复共轭与换叶分别作用为
\begin{equation}\label{eq:symmetry-domains}
z\mapsto\conj{z}:D_\pm\longrightarrow D_\mp,\qquad
h:z\mapsto-\dfrac{q_0^2}{z}:D_\pm\longrightarrow D_\mp.
\end{equation}

\paragraph{两种映射怎样作用于区域.}
共轭不改变 $|z|$, 而把 $\operatorname{Im}z$ 变为其相反数,
所以区域定义中的乘积满足
\[
\operatorname{Im}(\conj{z})
\left(1-\dfrac{q_0^2}{|\conj{z}|^2}\right)
=-\operatorname{Im}z\left(1-\dfrac{q_0^2}{|z|^2}\right).
\]
换叶映射则满足
\[
h(z)=-\dfrac{q_0^2\conj{z}}{|z|^2},
\qquad
\operatorname{Im}h(z)=\dfrac{q_0^2}{|z|^2}\operatorname{Im}z,
\qquad
|h(z)|^2=\dfrac{q_0^4}{|z|^2}.
\]
代回同一个区域判别式,
\begin{align*}
\operatorname{Im}h(z)
\left(1-\dfrac{q_0^2}{|h(z)|^2}\right)
&=\dfrac{q_0^2}{|z|^2}\operatorname{Im}z
\left(1-\dfrac{|z|^2}{q_0^2}\right)\\
&=-\operatorname{Im}z\left(1-\dfrac{q_0^2}{|z|^2}\right).
\end{align*}
所以两种映射分别都交换 $D_+$ 与 $D_-$.
它们的对易性可直接计算:
\[
h(\conj{z})=-\dfrac{q_0^2}{\conj{z}}
=\conj{\left(-\dfrac{q_0^2}{z}\right)}=\conj{h(z)}.
\]
连续作用两次换叶还给出
\[
h(h(z))=-\dfrac{q_0^2}{-\dfrac{q_0^2}{z}}=z.
\]
这些关系决定离散谱轨道, 也决定哪些列的等式能够在相同解析域中比较.

二者对易, 其复合保持 $D_\pm$. 以下先在普通连续谱上比较 Jost 归一化, 再将相应列关系延拓到解析域.

\begin{proposition}[Jost 解的共轭与换叶约化]\label{prop:symmetry}
令
\[
\sigma_2=\begin{pmatrix}0&-\ii\\\ii&0\end{pmatrix},\qquad
C(z)=\dfrac{\ii z}{q_0}\sigma_1.
\]
在定理 \ref{thm:jost} 的条件下,
\begin{equation}\label{eq:conjugation}
\chi_\pm(z)=\sigma_2\conj{\chi_\pm(\conj{z})}\sigma_2,\qquad
\mu_\pm(z)=\sigma_2\conj{\mu_\pm(\conj{z})}\sigma_2,
\end{equation}
\begin{equation}\label{eq:musheet}
\chi_\pm(h(z))=\chi_\pm(z)C(z),\qquad
\mu_\pm(h(z))=\mu_\pm(z)C(z).
\end{equation}
逐列写为
\begin{alignat}{2}
&\chi_{\pm,1}(z)=\ii\sigma_2\conj{\chi_{\pm,2}(\conj{z})},&\qquad
&\chi_{\pm,2}(z)=-\ii\sigma_2\conj{\chi_{\pm,1}(\conj{z})},\label{eq:jost-conjugate-columns}\\
&\chi_{\pm,1}(h(z))=\dfrac{\ii z}{q_0}\chi_{\pm,2}(z),&\qquad
&\chi_{\pm,2}(h(z))=\dfrac{\ii z}{q_0}\chi_{\pm,1}(z).\label{eq:jost-sheet-columns}
\end{alignat}
修正列 $\mu_{\pm,j}$ 满足相同关系.
\end{proposition}
\begin{proof}

\par\medskip\noindent\textbf{共轭约化中的矩阵乘法.}\quad
先计算三个常矩阵的基本关系:
\[
\sigma_2^2
=\begin{pmatrix}(-\ii)\ii&0\\0&\ii(-\ii)\end{pmatrix}=I,
\qquad
\sigma_2\sg
=\begin{pmatrix}0&\ii\\\ii&0\end{pmatrix},
\]
\[
\sigma_2\sg\sigma_2
=\begin{pmatrix}-1&0\\0&1\end{pmatrix}=-\sg,
\qquad
\sigma_2\sigma_1\sigma_2=-\sigma_1.
\]
对势矩阵则有
\begin{align*}
\sigma_2\conj{Q}
&=\begin{pmatrix}0&-\ii\\\ii&0\end{pmatrix}
\begin{pmatrix}0&\conj{q}\\-q&0\end{pmatrix}
=\begin{pmatrix}\ii q&0\\0&\ii\conj{q}\end{pmatrix},\\
\sigma_2\conj{Q}\sigma_2
&=\begin{pmatrix}0&q\\-\conj{q}&0\end{pmatrix}=Q.
\end{align*}
单值化有理函数的系数为实数, 所以
\[
\conj{k(\conj{z})}=k(z),
\qquad
\conj{\lambda(\conj{z})}=\lambda(z),
\qquad
\conj{F(k(\conj{z}))}=F(k(z)).
\]
于是
\begin{align*}
\sigma_2\conj{U(\conj{z})}\sigma_2
&=\sigma_2(\ii k(z)\sg+\conj{Q})\sigma_2\\
&=-\ii k(z)\sg+Q=U(z),\\
\sigma_2\conj{Y(\conj{z})}\sigma_2
&=\sigma_2\left(I+\dfrac{\ii q_0}{z}\sigma_1\right)\sigma_2\\
&=I-\dfrac{\ii q_0}{z}\sigma_1=Y(z).
\end{align*}
相位满足
\[
\conj{(-\lambda(\conj{z})x)}
=-\conj{\lambda(\conj{z})}
\bigl[x-\conj{F(k(\conj{z}))}t\bigr]
=-\lambda(z)[x-F(k(z))t]=(-\lambda(z)x).
\]
这里 $x,t$ 为实变量, 这一点在取共轭时也被使用了.

由 $\sigma_2\sg\sigma_2=-\sg$、$\sigma_2\conj{Q}\sigma_2=Q$ 及 $k(\conj{z})=\conj{k(z)}$,
\[
U(z)=\sigma_2\conj{U(\conj{z})}\sigma_2,\qquad
Y(z)=\sigma_2\conj{Y(\conj{z})}\sigma_2.
\]
又有 $\lambda(\conj{z})=\conj{\lambda(z)}$, 故
\[
\sigma_2\conj{Y(\conj{z})e^{-\ii\lambda(\conj{z})x\sg}}\sigma_2
=Y(z)e^{-\ii\lambda(z)x\sg}.
\]

\par\medskip\noindent\textbf{共轭指数的负号与列交换的负号.}\quad
取共轭先将 $\ii$ 变为 $-\ii$, 故
\[
\conj{e^{-\ii\lambda(\conj{z})x\sg}}
=e^{\ii\lambda(z)x\sg}.
\]
再与 $\sigma_2$ 共轭, 对角元交换:
\begin{align*}
\sigma_2 e^{\ii\lambda x\sg}\sigma_2
&=\sigma_2\begin{pmatrix}e^{\ii\lambda x}&0\\0&e^{-\ii\lambda x}\end{pmatrix}\sigma_2\\
&=\begin{pmatrix}e^{-\ii\lambda x}&0\\0&e^{\ii\lambda x}\end{pmatrix}
=e^{-\ii\lambda x\sg}.
\end{align*}
因为 $\sigma_2^2=I$, 乘积的约化可在中间插入 $I$:
\begin{align*}
\sigma_2\conj{Y(\conj{z})e^{-\ii\lambda(\conj{z})x\sg}}\sigma_2
&=(\sigma_2\conj{Y(\conj{z})}\sigma_2)
(\sigma_2 e^{\ii\lambda(z)x\sg}\sigma_2)\\
&=Y(z)e^{-\ii\lambda(z)x\sg}.
\end{align*}
对变换后的解求导,
\begin{align*}
\partial_x\bigl(\sigma_2\conj{\chi_\pm(\conj{z})}\sigma_2\bigr)
&=\sigma_2\conj{U(\conj{z})}\conj{\chi_\pm(\conj{z})}\sigma_2\\
&=U(z)\sigma_2\conj{\chi_\pm(\conj{z})}\sigma_2.
\end{align*}
这就是同一空间方程; 上式又给出同一端态,
所以可应用 Jost 唯一性.

共轭后的 Jost 解满足同一空间方程和同一归一化, 定理 \ref{thm:jost} 的唯一性给出 \eqref{eq:conjugation}.


\par\medskip\noindent\textbf{时间矩阵约化中有理项的核对.}\quad
对介质矩阵逐元计算,
\begin{align*}
\sigma_2\conj{R}\sigma_2
&=\sigma_2
\begin{pmatrix}\eta&-\conj{p}\\-p&-\eta\end{pmatrix}\sigma_2\\
&=\begin{pmatrix}-\eta&p\\\conj{p}&\eta\end{pmatrix}=-R.
\end{align*}
因此有理项经过约化时出现两个负号:
\[
\sigma_2\conj{\left(\dfrac{\ii R}{k(\conj{z})-a}\right)}\sigma_2
=-\dfrac{\ii\sigma_2\conj{R}\sigma_2}{k(z)-a}
=\dfrac{\ii R}{k(z)-a}.
\]
$\dfrac{\ii\Omega}{2}\sg$ 也由
$\conj{\ii}=-\ii$ 与 $\sigma_2\sg\sigma_2=-\sg$
的两个负号保持不变.
对多项式部分, 每个 $V_j$ 都是迹零矩阵,
其对角元纯虚且 $(V_j)_{21}=-\conj{(V_j)_{12}}$.
直接使用一般二阶乘法即有
\[
\sigma_2\conj{V_j}\sigma_2
=\begin{pmatrix}
\conj{(V_j)_{22}}&-\conj{(V_j)_{21}}\\
-\conj{(V_j)_{12}}&\conj{(V_j)_{11}}
\end{pmatrix}=V_j.
\]
例如 $q\conj{q_x}-q_x\conj{q}$ 的共轭是其相反数,
所以 $V_0$ 的对角元也确实纯虚.
这就逐项给出了 $V$ 的共轭约化.

换叶保持 $k$ 而使 $\lambda$ 变号. 直接由 \eqref{eq:E} 计算,
\begin{equation}\label{eq:Esheet}
Y(h(z))=I+\dfrac{\ii z}{q_0}\sigma_1=Y(z)C(z),\qquad
C(z)\sg=-\sg C(z).
\end{equation}

\par\medskip\noindent\textbf{换叶所需的全部代数恒等式.}\quad
把 $h(z)$ 代入 $k,\lambda$ 的有理表达式,
\begin{align*}
k(h(z))
&=\dfrac{1}{2}\left(-\dfrac{q_0^2}{z}
-\dfrac{q_0^2}{-\dfrac{q_0^2}{z}}\right)
=\dfrac{1}{2}\left(z-\dfrac{q_0^2}{z}\right)=k(z),\\
\lambda(h(z))
&=\dfrac{1}{2}\left(-\dfrac{q_0^2}{z}
+\dfrac{q_0^2}{-\dfrac{q_0^2}{z}}\right)
=-\dfrac{1}{2}\left(z+\dfrac{q_0^2}{z}\right)=-\lambda(z).
\end{align*}
因为 $F$ 只依赖 $k$, 可得 $F(k(h(z)))=F(k(z))$,
故 $(-\lambda(h(z))x)=-(-\lambda(z)x)$.
$Y$ 与 $C$ 的乘法也可直接展开:
\begin{align*}
Y(z)C(z)
&=\left(I-\dfrac{\ii q_0}{z}\sigma_1\right)
\dfrac{\ii z}{q_0}\sigma_1\\
&=\dfrac{\ii z}{q_0}\sigma_1
-\dfrac{\ii q_0}{z}\dfrac{\ii z}{q_0}\sigma_1^2\\
&=I+\dfrac{\ii z}{q_0}\sigma_1=Y(h(z)).
\end{align*}
其逆与两次换叶的一致性分别为
\[
C(z)^{-1}=-\dfrac{\ii q_0}{z}\sigma_1,
\qquad C(z)C(h(z))=I.
\]
由于
\[
\sigma_1\sg=\begin{pmatrix}0&-1\\1&0\end{pmatrix},
\qquad
\sg\sigma_1=\begin{pmatrix}0&1\\-1&0\end{pmatrix},
\]
有 $C\sg=-\sg C$.
把指数写为对角矩阵可逐元确认
\[
e^{-\ii\lambda x\sg}C
=C e^{\ii\lambda x\sg},
\qquad
C e^{-\ii\lambda x\sg}=e^{\ii\lambda x\sg}C.
\]
因此换叶不是简单地把指数换号后保持列号;
右乘 $C$ 同时交换了两列.

于是 $\chi_\pm(z)C(z)$ 在 $h(z)$ 处满足同一空间方程, 且
\[
Y(z)e^{-\ii\lambda(z)x\sg}C(z)
=Y(h(z))e^{-\ii\lambda(h(z))x\sg}.
\]
再次由唯一性得到 $\chi_\pm$ 的换叶关系. 去除指数时使用

\[
C(z)e^{-\ii\lambda(z)x\sg}=e^{\ii\lambda(z)x\sg}C(z)
\]
,
便得到 $\mu_\pm$ 的关系.
\par\medskip\noindent\textbf{修正 Jost 关系及四条列关系逐项展开.}\quad
从 $\chi(h(z))=\chi(z)C(z)$ 出发,
\begin{align*}
\mu(h(z))
&=\chi(h(z))e^{\ii\lambda(h(z))x\sg}\\
&=\chi(z)C(z)e^{-\ii\lambda(z)x\sg}\\
&=\chi(z)e^{\ii\lambda(z)x\sg}C(z)=\mu(z)C(z).
\end{align*}
共轭关系对 $\mu$ 的证明也保留两个指数负号:
\begin{align*}
\sigma_2\conj{\mu(\conj{z})}\sigma_2
&=\sigma_2\conj{\chi(\conj{z})}
 e^{-\ii\lambda(z)x\sg}\sigma_2\\
&=(\sigma_2\conj{\chi(\conj{z})}\sigma_2)
 (\sigma_2 e^{-\ii\lambda(z)x\sg}\sigma_2)\\
&=\chi(z)e^{\ii\lambda(z)x\sg}=\mu(z).
\end{align*}
把右乘 $\sigma_2$ 的结果写成两列,
\[
\conj{\chi(\conj{z})}\sigma_2
=\bigl(\ii\conj{\chi_2(\conj{z})},
-\ii\conj{\chi_1(\conj{z})}\bigr).
\]
再左乘 $\sigma_2$, 得
\[
\chi_1(z)=\ii\sigma_2\conj{\chi_2(\conj{z})},
\qquad
\chi_2(z)=-\ii\sigma_2\conj{\chi_1(\conj{z})}.
\]
右乘 $C$ 则没有两个不同的符号:
\[
(\chi_1,\chi_2)C
=\left(\dfrac{\ii z}{q_0}\chi_2,
\dfrac{\ii z}{q_0}\chi_1\right).
\]
相同乘法适用于 $\mu$.
对解析列比较时, 例如
$\mu_{-,1}(z)$ 与 $\mu_{-,2}(\conj{z})$,
右边的第二列在 $D_-$ 解析, 而共轭后的自变量属于 $D_+$;
所以两边确实在同一个区域上有意义.

右乘矩阵的两列分别给出 \eqref{eq:jost-conjugate-columns}--\eqref{eq:jost-sheet-columns}. 变换后的列具有 \eqref{eq:symmetry-domains} 所确定的解析域, 因此同一 Volterra 唯一性论证在各解析域内成立.
\end{proof}

\begin{proposition}[散射矩阵与反射数据的约化]\label{prop:scattering-symmetry}
散射矩阵满足
\begin{equation}\label{eq:scattering-symmetry}
S(z)=\sigma_2\conj{S(\conj{z})}\sigma_2,\qquad
S(h(z))=\sigma_1S(z)\sigma_1.
\end{equation}
特别地,
\begin{equation}\label{eq:sconj}
s_{22}(z)=\conj{s_{11}(\conj{z})},\qquad
s_{12}(z)=-\conj{s_{21}(\conj{z})},
\end{equation}
\begin{equation}\label{eq:ssheet}
s_{22}(h(z))=s_{11}(z),\qquad s_{12}(h(z))=s_{21}(z),
\end{equation}
\begin{equation}\label{eq:reflection-symmetry}
\rho(z)=-\conj{\widetilde\rho(\conj{z})},\qquad
\rho(h(z))=\widetilde\rho(z).
\end{equation}
透射系数的关系在相应解析域内成立. 在实谱的普通点,
\begin{equation}\label{eq:scattering-unitarity}
S(z)^\dagger S(z)=I,\qquad
|s_{11}(z)|^2+|s_{21}(z)|^2=1.
\end{equation}
\end{proposition}
\begin{proof}

\par\medskip\noindent\textbf{由基底的约化得到散射矩阵的约化.}\quad
先在 $\conj{z}$ 处写散射关系并取共轭,
\[
\conj{\chi_+(\conj{z})}
=\conj{\chi_-(\conj{z})}\conj{S(\conj{z})}.
\]
左右乘 $\sigma_2$, 并在乘积中间插入
$\sigma_2^2=I$, 得
\begin{align*}
\chi_+(z)
&=(\sigma_2\conj{\chi_-(\conj{z})}\sigma_2)
(\sigma_2\conj{S(\conj{z})}\sigma_2)\\
&=\chi_-(z)(\sigma_2\conj{S(\conj{z})}\sigma_2).
\end{align*}
与 $\chi_+=\chi_-S$ 比较, 左乘可逆的 $\chi_-^{-1}$,
得到共轭约化.
换叶时在 $h(z)$ 处写散射关系,
\[
\chi_+(z)C(z)=\chi_-(z)C(z)S(h(z)).
\]
同时原关系右乘 $C(z)$ 给出
$\chi_+(z)C(z)=\chi_-(z)S(z)C(z)$.
两式消去 $\chi_-$ 后,
\[
C(z)S(h(z))=S(z)C(z),
\qquad
S(h(z))=C(z)^{-1}S(z)C(z).
\]
$C$ 中非零的标量因子在共轭中抵消, 因而只留下
$\sigma_1S\sigma_1$.

将 \eqref{eq:conjugation} 代入 $\chi_+=\chi_-S$, 并利用 $\chi_-$ 的可逆性, 得到 \eqref{eq:scattering-symmetry} 的第一式. 比较
\[
\chi_+(z)C(z)=\chi_-(z)S(z)C(z)
=\chi_-(z)C(z)S(h(z)),
\]
得到
\[
S(h(z))=C(z)^{-1}S(z)C(z)=\sigma_1S(z)\sigma_1.
\]

\par\medskip\noindent\textbf{把两个矩阵约化完全展开为四个矩阵元.}\quad
共轭约化的右侧实际为
\begin{align*}
\sigma_2\conj{S(\conj{z})}\sigma_2
&=\begin{pmatrix}
\conj{s_{22}(\conj{z})}&-\conj{s_{21}(\conj{z})}\\
-\conj{s_{12}(\conj{z})}&\conj{s_{11}(\conj{z})}
\end{pmatrix}.
\end{align*}
故除了原文列出的两式之外, 另外两式为
\[
s_{11}(z)=\conj{s_{22}(\conj{z})},
\qquad
s_{21}(z)=-\conj{s_{12}(\conj{z})}.
\]
换叶约化则给出
\[
\sigma_1S(z)\sigma_1
=\begin{pmatrix}s_{22}(z)&s_{21}(z)\\s_{12}(z)&s_{11}(z)\end{pmatrix}.
\]
所以四个矩阵元分别是
\begin{align*}
s_{11}(h(z))&=s_{22}(z),&s_{12}(h(z))&=s_{21}(z),\\
s_{21}(h(z))&=s_{12}(z),&s_{22}(h(z))&=s_{11}(z).
\end{align*}
按反射系数的分母逐项相除,
\begin{align*}
\rho(z)
&=\dfrac{-\conj{s_{12}(\conj{z})}}{\conj{s_{22}(\conj{z})}}
=-\conj{\left(\dfrac{s_{12}(\conj{z})}{s_{22}(\conj{z})}\right)}
=-\conj{\widetilde\rho(\conj{z})},\\
\rho(h(z))
&=\dfrac{s_{21}(h(z))}{s_{11}(h(z))}
=\dfrac{s_{12}(z)}{s_{22}(z)}=\widetilde\rho(z).
\end{align*}
负号来自 $\sigma_2$ 的非对角元, 不来自行列式约定的改变.
在非实圆谱上应保留 $\conj{z}$ 这个自变量;
只有 $z$ 为实时才可把它替换为 $z$.

逐元展开即得 \eqref{eq:sconj}--\eqref{eq:ssheet}; 除以相应透射系数得到 \eqref{eq:reflection-symmetry}. 解析透射系数的关系也可直接对 \eqref{eq:wronskian} 应用 Jost 列的约化得到.


\par\medskip\noindent\textbf{实轴上的酉性怎样逐项得到.}\quad
在实 $z$ 上,
\[
U=\begin{pmatrix}-\ii k&q\\-\conj{q}&\ii k\end{pmatrix},
\qquad
U^\dagger=\begin{pmatrix}\ii k&-q\\\conj{q}&-\ii k\end{pmatrix}=-U.
\]
所以
\begin{align*}
\partial_x(\chi_\pm^\dagger\chi_\pm)
&=(U\chi_\pm)^\dagger\chi_\pm
+\chi_\pm^\dagger U\chi_\pm\\
&=\chi_\pm^\dagger U^\dagger\chi_\pm
+\chi_\pm^\dagger U\chi_\pm=0.
\end{align*}
背景矩阵的乘积为
\begin{align*}
Y^\dagger Y
&=\begin{pmatrix}1&\dfrac{\ii q_0}{z}\\
\dfrac{\ii q_0}{z}&1\end{pmatrix}
\begin{pmatrix}1&-\dfrac{\ii q_0}{z}\\
-\dfrac{\ii q_0}{z}&1\end{pmatrix}\\
&=\begin{pmatrix}1+\dfrac{q_0^2}{z^2}&0\\
0&1+\dfrac{q_0^2}{z^2}\end{pmatrix}=\gamma I.
\end{align*}
由于 $(-\lambda x)$ 为实数, $e^{-\ii\lambda x\sg}$ 是酉矩阵,
在归一化端
\[
\chi_\pm^\dagger\chi_\pm
\longrightarrow e^{\ii\lambda x\sg}\gamma I e^{-\ii\lambda x\sg}=\gamma I.
\]
空间守恒性使该等式对所有有限 $x$ 成立.
代入 $\chi_+=\chi_-S$,
\[
\gamma I=\chi_+^\dagger\chi_+
=S^\dagger\chi_-^\dagger\chi_-S
=\gamma S^\dagger S.
\]
$\gamma=1+\dfrac{q_0^2}{z^2}>0$, 故可除以它.
最后逐元展开得到
\[
S^\dagger S
=\begin{pmatrix}
|s_{11}|^2+|s_{21}|^2&\conj{s_{11}}s_{12}+\conj{s_{21}}s_{22}\\
\conj{s_{12}}s_{11}+\conj{s_{22}}s_{21}&|s_{12}|^2+|s_{22}|^2
\end{pmatrix}=I.
\]
左上元就是原文的标量恒等式, 右上元则给出两列的正交关系.
结合 $\widetilde\rho=-\conj\rho$ 还可复核
\[
|s_{11}|^2(1+|\rho|^2)=1,
\]
这与行列式关系的实轴取值一致.

在实 $z$ 上, $k,\lambda$ 为实数, 且 $U^\dagger=-U$. 因而
\[
\partial_x(\chi_\pm^\dagger\chi_\pm)
=\chi_\pm^\dagger(U^\dagger+U)\chi_\pm=0,\qquad
Y^\dagger Y=\gamma I.
\]
振荡归一化给出 $\chi_\pm^\dagger\chi_\pm=\gamma I$. 代入散射关系并除以 $\gamma>0$, 即得 \eqref{eq:scattering-unitarity}.
\end{proof}
\subsection{离散谱、范数数据与留数}\label{ch:discrete}
固定 $t$, 假设透射系数的离散零点有限且均为简单零点.离散谱指空间方程在 $\mathbb C\mathbin{\backslash}\Sigma$ 中存在非零 $L^2(\mathbb R)$ 解的谱点集合.

\begin{proposition}[四点离散谱轨道]\label{prop:quartet}
离散谱恰为 $s_{22}$ 在 $D_+$ 中及 $s_{11}$ 在 $D_-$ 中的零点. 若 $\zeta\in D_+$ 是 $s_{22}$ 的简单零点, 则
\begin{equation}\label{eq:quartet}
\left\{\zeta,\conj{\zeta},-\dfrac{q_0^2}{\zeta},-\dfrac{q_0^2}{\conj{\zeta}}\right\}
\end{equation}
构成四点互异的完整轨道. 其中 $\zeta$ 与 $h(\conj{\zeta})$ 属于 $D_+$, 其余两点属于 $D_-$. 每个轨道选取外部上半平面的代表 $\zeta_j$, $j=1,\ldots,N$, 再令
\begin{equation}\label{eq:discrete-indexing}
\zeta_{N+j}=h(\conj{\zeta_j}),\qquad
Z=\{\zeta_j,\conj{\zeta_j}:j=1,\ldots,2N\}.
\end{equation}
每个 $D_+$ 零点满足非零列相关关系
\begin{equation}\label{eq:norming-definition}
\chi_{+,2}(\zeta_j)=b_j(t)\chi_{-,1}(\zeta_j),\qquad
c_j(t)=\dfrac{b_j(t)}{s_{22}'(\zeta_j)},\qquad
C_j(x,t)=c_j(t)e^{-2\ii\lambda(\zeta_j)x}.
\end{equation}
这里 $j=1,\ldots,2N$, $b_j(t),c_j(t)$ 与 $x$ 无关, 撇号表示对 $z$ 求导.
\end{proposition}
\begin{proof}
在 $D_+$ 中, $\operatorname{Im}\lambda>0$, 因而
\[
\chi_{-,1}\sim Y_1e^{-\ii\lambda x}\quad(x\to-\infty),\qquad
\chi_{+,2}\sim Y_2e^{\ii\lambda x}\quad(x\to+\infty)
\]
分别指数衰减. 任意另一个半轴 $L^2$ 解与该衰减列的 Wronskian 为常数; Cauchy--Schwarz 不等式又保证该 Wronskian 在半轴上可积, 因而此常数只能为零. 两端的 $L^2$ 解空间因此各为一维, 一个解在两端均属 $L^2$ 当且仅当这两列线性相关. 由 \eqref{eq:wronskian}, 这又等价于 $s_{22}(\zeta)=0$. 在 $D_-$ 中交换两列即得 $s_{11}$ 的结论.

共轭与换叶各自交换两个解析域, 二者复合保留解析域. 由 \eqref{eq:sconj}--\eqref{eq:ssheet},
\[
s_{22}(\zeta)=0
\quad\Longrightarrow\quad
s_{11}(\conj{\zeta})=s_{11}(h(\zeta))=s_{22}(h(\conj{\zeta}))=0.
\]
$h'(z)=\dfrac{q_0^2}{z^2}\ne0$ 保持零点的简单性. 共轭的不动点位于实轴, 换叶的不动点为 $\pm\ii q_0$, 而 $h(\conj{z})=z$ 将导致 $|z|^2=-q_0^2$. 因而 $\Sigma$ 外的四点互异, 且每个轨道恰有一个外部上半平面的代表.


Jost 列是非零空间解, 所以 $b_j(t)\ne0$.
对 $\chi_{+,2}=b_j(t)\chi_{-,1}$ 求 $x$ 导数并使用同一空间方程,
得到 $(b_j)_x\chi_{-,1}=0$. 因而 $b_j$ 与 $x$ 无关.
固定 $t$ 时, 下文可省略 $b_j(t),c_j(t)$ 的时间自变量.
它们的时间演化将在第 \ref{ch:scattering} 节由方程证明.

\end{proof}

\begin{proposition}[简单零点的留数及常数约化]\label{prop:respair}
在式 \eqref{eq:norming-definition} 的条件下, 相应解析商满足
\begin{align}
\Res_{z=\zeta_j}\dfrac{\mu_{+,2}(z)}{s_{22}(z)}
&=C_j\mu_{-,1}(\zeta_j),\label{eq:resplus}\\
\Res_{z=\conj{\zeta_j}}\dfrac{\mu_{+,1}(z)}{s_{11}(z)}
&=-\conj{C_j}\mu_{-,2}(\conj{\zeta_j}).\label{eq:resminus}
\end{align}
若 $\widehat\zeta=h(\conj{\zeta})$ 为同一轨道的另一个 $D_+$ 点, 则
\begin{equation}\label{eq:cpair}
b_{\widehat\zeta}=-\conj{b_\zeta},\qquad
c_{\widehat\zeta}=-\dfrac{q_0^2}{(\conj{\zeta})^2}\conj{c_\zeta}.
\end{equation}
相应时空留数系数满足
\begin{equation}\label{eq:Cpair}
C_{\widehat\zeta}=-\dfrac{q_0^2}{(\conj{\zeta})^2}\conj{C_\zeta}.
\end{equation}
\end{proposition}
\begin{proof}
由 $\chi_{+,2}=b_j\chi_{-,1}$ 及 $\chi_1=\mu_1e^{-\ii\lambda x}$、$\chi_2=\mu_2e^{\ii\lambda x}$, 得到
\[
\mu_{+,2}(\zeta_j)=b_je^{-2\ii\lambda(\zeta_j)x}\mu_{-,1}(\zeta_j).
\]
在简单零点处展开,
\[
s_{22}(z)=s_{22}'(\zeta_j)(z-\zeta_j)+O((z-\zeta_j)^2),\qquad
s_{22}'(\zeta_j)\ne0.
\]
因此
\[
\Res_{z=\zeta_j}\dfrac{\mu_{+,2}(z)}{s_{22}(z)}
=\lim\limits_{z\to\zeta_j}\dfrac{(z-\zeta_j)\mu_{+,2}(z)}{s_{22}(z)}
=\dfrac{\mu_{+,2}(\zeta_j)}{s_{22}'(\zeta_j)},
\]
代入列相关关系即得 \eqref{eq:resplus}. 对列相关关系应用 \eqref{eq:jost-conjugate-columns}, 在 $\conj{\zeta_j}$ 处得到
\[
\chi_{+,1}=-\conj{b_j}\chi_{-,2},\qquad
s_{11}'(\conj{\zeta_j})=\conj{s_{22}'(\zeta_j)}.
\]
这里的导数关系也可由上述 Taylor 展开取共轭直接得到. 又因 $(-\lambda(\conj{\zeta_j})x)=\conj{(-\lambda(\zeta_j)x)}$,
\[
\dfrac{\mu_{+,1}(\conj{\zeta_j})}{s_{11}'(\conj{\zeta_j})}
=-\conj{c_j}e^{2\ii\lambda(\conj{\zeta_j})x}
\mu_{-,2}(\conj{\zeta_j})
=-\conj{C_j}\mu_{-,2}(\conj{\zeta_j}),
\]
即为 \eqref{eq:resminus}.

再在 $z=\conj{\zeta}$ 使用换叶关系. 它给出
\[
\chi_{+,2}(\widehat\zeta)
=\dfrac{\ii\conj{\zeta}}{q_0}\chi_{+,1}(\conj{\zeta})
=-\conj{b_\zeta}\dfrac{\ii\conj{\zeta}}{q_0}\chi_{-,2}(\conj{\zeta})
=-\conj{b_\zeta}\chi_{-,1}(\widehat\zeta).
\]
对恒等式 $s_{22}(h(z))=s_{11}(z)$ 求导, 有
\[
s_{22}'(\widehat\zeta)h'(\conj{\zeta})=\conj{s_{22}'(\zeta)},\qquad
h'(\conj{\zeta})=\dfrac{q_0^2}{(\conj{\zeta})^2}.
\]
链式因子也可从两次展开直接看出:
\[
\begin{aligned}
h(z)-\widehat\zeta
&=h'(\conj{\zeta})(z-\conj{\zeta})+O((z-\conj{\zeta})^2),\\
s_{22}(h(z))
&=s_{22}'(\widehat\zeta)h'(\conj{\zeta})(z-\conj{\zeta})+O((z-\conj{\zeta})^2).
\end{aligned}
\]
因此
\[
c_{\widehat\zeta}
=\dfrac{-\conj{b_\zeta}}{s_{22}'(\widehat\zeta)}
=-h'(\conj{\zeta})\conj{c_\zeta},
\]
即为 \eqref{eq:cpair}. 最后 $\lambda(\widehat\zeta)=-\conj{\lambda(\zeta)}$, 故
\[
C_{\widehat\zeta}
=-\dfrac{q_0^2}{(\conj{\zeta})^2}\conj{c_\zeta}
e^{2\ii\conj{\lambda(\zeta)}x}
=-\dfrac{q_0^2}{(\conj{\zeta})^2}\conj{C_\zeta},
\]
得到 \eqref{eq:Cpair}. 对全部 $D_+$ 零点应用 \eqref{eq:resplus}--\eqref{eq:resminus}, 即给出完整四点轨道的留数条件.
\end{proof}


\subsubsection*{本部分全部谱点与留数公式的逐式计算}
\paragraph{式 \eqref{eq:quartet} 之前: 两列的指数衰减.}
固定 $t$ 与谱点, 原文相位直接给出
\begin{align*}
e^{-\ii\lambda x}&=e^{-\ii\lambda x}e^{\ii\lambda Ft},&
|e^{-\ii\lambda x}|&=e^{x\operatorname{Im}\lambda}
 e^{-t\operatorname{Im}(\lambda F)},\\
e^{\ii\lambda x}&=e^{\ii\lambda x}e^{-\ii\lambda Ft},&
|e^{\ii\lambda x}|&=e^{-x\operatorname{Im}\lambda}
 e^{t\operatorname{Im}(\lambda F)}.
\end{align*}
在 $D_+$ 中 $\operatorname{Im}\lambda>0$.
第一行随 $x\to-\infty$ 衰减, 第二行随 $x\to+\infty$ 衰减.
对应的 $\mu$ 列在归一化半轴有界, 因而
\begin{align*}
\int_{-\infty}^0\|\chi_{-,1}(x)\|_1^2\,\dd x
&\le\|\mu_{-,1}\|_{L^\infty(-\infty,0)}^2
 e^{-2t\operatorname{Im}(\lambda F)}
 \int_{-\infty}^0 e^{2x\operatorname{Im}\lambda}\,\dd x\\
&=\dfrac{\|\mu_{-,1}\|_{L^\infty(-\infty,0)}^2
 e^{-2t\operatorname{Im}(\lambda F)}}{2\operatorname{Im}\lambda}<\infty,\\
\int_0^{+\infty}\|\chi_{+,2}(x)\|_1^2\,\dd x
&\le\|\mu_{+,2}\|_{L^\infty(0,+\infty)}^2
 e^{2t\operatorname{Im}(\lambda F)}
 \int_0^{+\infty} e^{-2x\operatorname{Im}\lambda}\,\dd x\\
&=\dfrac{\|\mu_{+,2}\|_{L^\infty(0,+\infty)}^2
 e^{2t\operatorname{Im}(\lambda F)}}{2\operatorname{Im}\lambda}<\infty.
\end{align*}
这里把半轴在 $0$ 分开只是选择积分区间, 没有对变量换元.
两个指数积分分别由
\[
\left.\dfrac{e^{2x\operatorname{Im}\lambda}}{2\operatorname{Im}\lambda}
\right|_{-\infty}^0
=\dfrac{1}{2\operatorname{Im}\lambda},
\qquad
\left.-\dfrac{e^{-2x\operatorname{Im}\lambda}}{2\operatorname{Im}\lambda}
\right|_0^{+\infty}
=\dfrac{1}{2\operatorname{Im}\lambda}
\]
计算得到.
在 $D_-$ 中 $\operatorname{Im}\lambda<0$,
$\chi_{+,1}$ 在正端平方可积, $\chi_{-,2}$ 在负端平方可积.
把上面的积分指数符号与端点同时反转,
两个积分的值都变为 $-\dfrac{1}{2\operatorname{Im}\lambda}>0$.

\paragraph{平方可积条件怎样等价于透射零点.}
第4节对行列式的求导计算适用于同一系统的任意两列,
因此任意两个空间解的 Wronskian 都是常数.
如果一个是衰减 Jost 列, 另一个也是同一半轴上的 $L^2$ 解,
二阶行列式绝对值不超过两列欧氏范数的乘积.
将该不等式在半轴上积分, 再用积分型 Cauchy--Schwarz 不等式,
得到的上界是两列各自平方积分的平方根之积, 因而有限.
这里没有对另一解另起辅助记号; 所用不等式就是行列式与范数的基本估计.
一个非零常数在无限长半轴上的绝对值积分为无穷,
所以这个 Wronskian 常数只能为零.
于是任意半轴 $L^2$ 解都与对应 Jost 衰减列线性相关.
两列在一处相关后, 空间线性初值问题的唯一性使同一个关系在全轴成立.

因此 $D_+$ 内存在全轴 $L^2$ 解, 当且仅当
\[
\det(\chi_{-,1},\chi_{+,2})=0.
\]
用原文 Wronskian 公式直接替换,
\[
\det(\chi_{-,1},\chi_{+,2})=\gamma s_{22},
\qquad \gamma\ne0\quad\hbox{在 }D_+,
\]
可知这与 $s_{22}=0$ 等价.
反方向也具体成立: 如果 $s_{22}=0$, 两个非零列相关,
$\chi_{-,1}$ 在负端衰减, 在正端又是 $\chi_{+,2}$ 的非零常数倍,
所以它在全轴平方可积.
在 $D_-$ 中用
\[
\det(\chi_{+,1},\chi_{-,2})=\gamma s_{11}
\]
得到同样结论.
没有把一般非实谱上的指数增长列误认为 $L^2$ 列.

\paragraph{式 \eqref{eq:quartet}: 四点的零点类型与区域.}
由共轭约化先有
\[
s_{11}(\conj{\zeta})=\conj{s_{22}(\zeta)}=0.
\]
再由换叶约化,
\[
s_{11}(h(\zeta))=s_{22}(\zeta)=0,
\qquad
s_{22}(h(\conj{\zeta}))=s_{11}(\conj{\zeta})=0.
\]
两种映射分别交换 $D_+$ 与 $D_-$,
所以四点的归属逐个是
\[
\zeta\in D_+,
\qquad \conj{\zeta}\in D_-,
\qquad h(\zeta)\in D_-,
\qquad h(\conj{\zeta})\in D_+.
\]
共轭两次与换叶两次都回到原点, 且两种变换对易,
所以再对这四点作任意一次约化不会产生第五个点.

可能的重合只有以下三类, 其余情形可再作一次约化化到这三类:
\begin{align*}
\zeta=\conj{\zeta}
&\ \Longrightarrow\ \operatorname{Im}\zeta=0,\\
\zeta=h(\zeta)
&\ \Longrightarrow\ \zeta=-\dfrac{q_0^2}{\zeta}
\ \Longrightarrow\ \zeta^2=-q_0^2
\ \Longrightarrow\ \zeta=\pm\ii q_0,\\
\zeta=h(\conj{\zeta})
&\ \Longrightarrow\ \zeta\conj{\zeta}=-q_0^2
\ \Longrightarrow\ |\zeta|^2=-q_0^2.
\end{align*}
第一类位于实连续谱上, 第二类位于分支点上,
第三类因左侧非负而右侧为负不可能成立.
所以连续谱外四点互异.

\paragraph{式 \eqref{eq:discrete-indexing}: 为什么每轨道只有一个所选代表.}
$D_+$ 的两块区域是外部上半平面与内部下半平面.
若 $\zeta$ 已位于外部上半平面, 它就是代表.
若 $\zeta$ 位于内部下半平面, 则
\[
|h(\conj{\zeta})|=\dfrac{q_0^2}{|\zeta|}>q_0,
\qquad
\operatorname{Im}h(\conj{\zeta})
=-\dfrac{q_0^2}{|\zeta|^2}\operatorname{Im}\zeta>0.
\]
所以同一轨道的另一个 $D_+$ 点恰好是外部上半平面代表.
另外两点在 $D_-$, 因而不可能再成为这种代表.
枚举 $N$ 个轨道后, 前 $N$ 个点是这些代表,
后 $N$ 个点正好由 $h(\conj{\zeta_j})$ 补齐,
再加上共轭点就得到 $4N$ 个完整离散谱点.

\paragraph{范数数据的空间常数性质.}

Jost 列是非零空间解, 所以 $b_j(t)\ne0$.
对 $\chi_{+,2}=b_j(t)\chi_{-,1}$ 求 $x$ 导数并使用同一空间方程,
得到 $(b_j)_x\chi_{-,1}=0$. 因而 $b_j$ 与 $x$ 无关.
固定 $t$ 时, 下文可省略 $b_j(t),c_j(t)$ 的时间自变量.
它们的时间演化将在第 \ref{ch:scattering} 节由方程证明.

\paragraph{式 \eqref{eq:resplus}: 指数转换与分子分母的 Taylor 计算.}
将原 Jost 列与修正列的定义逐个代入,
\[
\mu_{+,2}(\zeta_j)e^{\ii\lambda(\zeta_j)x}
=b_j\mu_{-,1}(\zeta_j)e^{-\ii\lambda(\zeta_j)x}.
\]
两侧乘 $e^{-\ii\lambda(\zeta_j)x}$ 后,
\[
\mu_{+,2}(\zeta_j)
=b_je^{-2\ii\lambda(\zeta_j)x}\mu_{-,1}(\zeta_j).
\]
现在同时展开分子与分母:
\begin{align*}
\mu_{+,2}(z)
&=\mu_{+,2}(\zeta_j)+(z-\zeta_j)\mu_{+,2}'(\zeta_j)
+O((z-\zeta_j)^2),\\
s_{22}(z)
&=(z-\zeta_j)\left(s_{22}'(\zeta_j)
+\dfrac{s_{22}''(\zeta_j)}{2}(z-\zeta_j)
+O((z-\zeta_j)^2)\right).
\end{align*}
由于 $s_{22}'(\zeta_j)\ne0$, 对括号求倒数并相乘, 得
\begin{align*}
\dfrac{1}{s_{22}(z)}
&=\dfrac{1}{s_{22}'(\zeta_j)(z-\zeta_j)}
-\dfrac{s_{22}''(\zeta_j)}{2s_{22}'(\zeta_j)^2}
+O(z-\zeta_j),\\
\dfrac{\mu_{+,2}(z)}{s_{22}(z)}
&=\dfrac{\mu_{+,2}(\zeta_j)}{s_{22}'(\zeta_j)(z-\zeta_j)}
+\dfrac{\mu_{+,2}'(\zeta_j)}{s_{22}'(\zeta_j)}\\
&\quad-\dfrac{\mu_{+,2}(\zeta_j)s_{22}''(\zeta_j)}{2s_{22}'(\zeta_j)^2}
+O(z-\zeta_j).
\end{align*}
倒数展开的一次系数也可以通过乘回验证:
常数项为 1, 一次项是
$\dfrac{s_{22}''}{2s_{22}'}-\dfrac{s_{22}''}{2s_{22}'}=0$.
因此极点系数正是
\[
\dfrac{\mu_{+,2}(\zeta_j)}{s_{22}'(\zeta_j)}
=\dfrac{b_j}{s_{22}'(\zeta_j)}e^{-2\ii\lambda(\zeta_j)x}
\mu_{-,1}(\zeta_j)=C_j\mu_{-,1}(\zeta_j).
\]
分子在该点非零, 所以这里不是可去奇点.

\paragraph{式 \eqref{eq:resminus}: 共轭产生的负号逐列检查.}
原相关式取共轭后为
\[
\conj{\chi_{+,2}(\zeta_j)}
=\conj{b_j}\conj{\chi_{-,1}(\zeta_j)}.
\]
左乘 $\ii\sigma_2$,
\begin{align*}
\chi_{+,1}(\conj{\zeta_j})
&=\conj{b_j}\ii\sigma_2\conj{\chi_{-,1}(\zeta_j)},\\
\chi_{-,2}(\conj{\zeta_j})
&=-\ii\sigma_2\conj{\chi_{-,1}(\zeta_j)}.
\end{align*}
第二行的 $-\ii$ 与第一行的 $+\ii$ 相反,
因此
\[
\chi_{+,1}(\conj{\zeta_j})
=-\conj{b_j}\chi_{-,2}(\conj{\zeta_j}).
\]
改写成修正列,
\[
\mu_{+,1}(\conj{\zeta_j})
=-\conj{b_j}e^{2\ii\lambda(\conj{\zeta_j})x}
\mu_{-,2}(\conj{\zeta_j}).
\]
现在被除的列是第一列, 对应的空间指数为 $e^{2\ii\lambda x}$.

\paragraph{共轭零点处的导数由 Taylor 系数比较得到.}
不能把非解析映射 $z\mapsto\conj{z}$ 当作解析函数套用复导数链式法则.
正确计算从
$s_{11}(z)=\conj{s_{22}(\conj{z})}$ 的局部展开开始:
\begin{align*}
s_{22}(\conj{z})
&=s_{22}'(\zeta_j)(\conj{z}-\zeta_j)
+O(|\conj{z}-\zeta_j|^2),\\
s_{11}(z)
&=\conj{s_{22}'(\zeta_j)}(z-\conj{\zeta_j})
+O(|z-\conj{\zeta_j}|^2).
\end{align*}
因此
\[
s_{11}'(\conj{\zeta_j})=\conj{s_{22}'(\zeta_j)}\ne0.
\]
又有
\begin{align*}
\conj{C_j}
&=\conj{c_j}\conj{e^{-2\ii\lambda(\zeta_j)x}}\\
&=\conj{c_j}e^{2\ii\conj{\lambda(\zeta_j)}x}
=\conj{c_j}e^{2\ii\lambda(\conj{\zeta_j})x}.
\end{align*}
所以第二个简单极点的留数为
\begin{align*}
\dfrac{\mu_{+,1}(\conj{\zeta_j})}{s_{11}'(\conj{\zeta_j})}
&=-\dfrac{\conj{b_j}}{\conj{s_{22}'(\zeta_j)}}
 e^{2\ii\lambda(\conj{\zeta_j})x}\mu_{-,2}(\conj{\zeta_j})\\
&=-\conj{C_j}\mu_{-,2}(\conj{\zeta_j}).
\end{align*}
共轭后的简单性也在非零的一次 Taylor 系数中同时得到证明.

\paragraph{式 \eqref{eq:cpair}: 先算列相关常数, 再算导数.}
在 $\conj{\zeta}$ 处使用两条换叶列关系,
\begin{align*}
\chi_{+,2}(\widehat\zeta)
&=\dfrac{\ii\conj{\zeta}}{q_0}\chi_{+,1}(\conj{\zeta}),\\
\chi_{-,1}(\widehat\zeta)
&=\dfrac{\ii\conj{\zeta}}{q_0}\chi_{-,2}(\conj{\zeta}).
\end{align*}
把共轭点的列相关关系代入第一行,
\begin{align*}
\chi_{+,2}(\widehat\zeta)
&=-\conj{b_\zeta}\dfrac{\ii\conj{\zeta}}{q_0}
\chi_{-,2}(\conj{\zeta})\\
&=-\conj{b_\zeta}\chi_{-,1}(\widehat\zeta).
\end{align*}
右侧列非零, 所以相关倍数唯一,
$b_{\widehat\zeta}=-\conj{b_\zeta}$.

换叶的差值可以精确通分:
\begin{align*}
h(z)-h(\conj{\zeta})
&=-\dfrac{q_0^2}{z}+\dfrac{q_0^2}{\conj{\zeta}}
=\dfrac{q_0^2(z-\conj{\zeta})}{z\conj{\zeta}},\\
\lim\limits_{z\to\conj{\zeta}}
\dfrac{h(z)-h(\conj{\zeta})}{z-\conj{\zeta}}
&=\dfrac{q_0^2}{(\conj{\zeta})^2}.
\end{align*}
对 $s_{22}(h(z))=s_{11}(z)$ 写差商,
\[
\dfrac{s_{22}(h(z))-s_{22}(\widehat\zeta)}{h(z)-\widehat\zeta}
\dfrac{h(z)-\widehat\zeta}{z-\conj{\zeta}}
=\dfrac{s_{11}(z)-s_{11}(\conj{\zeta})}{z-\conj{\zeta}}.
\]
取极限后,
\[
s_{22}'(\widehat\zeta)\dfrac{q_0^2}{(\conj{\zeta})^2}
=\conj{s_{22}'(\zeta)},
\qquad
s_{22}'(\widehat\zeta)
=\dfrac{(\conj{\zeta})^2}{q_0^2}\conj{s_{22}'(\zeta)}.
\]
两因子都不为零, 所以另一 $D_+$ 零点仍然简单.
把结果代入 $c$ 的定义,
\begin{align*}
c_{\widehat\zeta}
&=\dfrac{-\conj{b_\zeta}}
{\dfrac{(\conj{\zeta})^2}{q_0^2}\conj{s_{22}'(\zeta)}}\\
&=-\dfrac{q_0^2}{(\conj{\zeta})^2}
\conj{\left(\dfrac{b_\zeta}{s_{22}'(\zeta)}\right)}
=-\dfrac{q_0^2}{(\conj{\zeta})^2}\conj{c_\zeta}.
\end{align*}
这里分母导数乘逆因子, 但再除以导数使留数系数乘正向因子.
所有撇号都对原文的 $z$ 求导, 没有改用 $k$ 导数.

\paragraph{式 \eqref{eq:Cpair}: 相位与四个留数的完整检查.}
由 $k$ 在换叶下不变而 $\lambda$ 变号,
\[
(-\lambda(\widehat\zeta)x)=-\conj{(-\lambda(\zeta)x)},
\qquad
e^{-2\ii\lambda(\widehat\zeta)x}
=e^{2\ii\conj{\lambda(\zeta)}x}
=\conj{e^{-2\ii\lambda(\zeta)x}}.
\]
所以
\begin{align*}
C_{\widehat\zeta}
&=c_{\widehat\zeta}e^{-2\ii\lambda(\widehat\zeta)x}\\
&=-\dfrac{q_0^2}{(\conj{\zeta})^2}\conj{c_\zeta}
\conj{e^{-2\ii\lambda(\zeta)x}}
=-\dfrac{q_0^2}{(\conj{\zeta})^2}\conj{C_\zeta}.
\end{align*}
完整四点留数逐一为
\begin{align*}
\Res_{z=\zeta}\dfrac{\mu_{+,2}(z)}{s_{22}(z)}
&=C_\zeta\mu_{-,1}(\zeta),\\
\Res_{z=\conj{\zeta}}\dfrac{\mu_{+,1}(z)}{s_{11}(z)}
&=-\conj{C_\zeta}\mu_{-,2}(\conj{\zeta}),\\
\Res_{z=h(\conj{\zeta})}\dfrac{\mu_{+,2}(z)}{s_{22}(z)}
&=-\dfrac{q_0^2}{(\conj{\zeta})^2}\conj{C_\zeta}
\mu_{-,1}(h(\conj{\zeta})),\\
\Res_{z=h(\zeta)}\dfrac{\mu_{+,1}(z)}{s_{11}(z)}
&=\dfrac{q_0^2}{\zeta^2}C_\zeta\mu_{-,2}(h(\zeta)).
\end{align*}
最后一行的系数来自
\[
-\conj{C_{\widehat\zeta}}
=-\conj{\left(-\dfrac{q_0^2}{(\conj{\zeta})^2}\conj{C_\zeta}\right)}
=\dfrac{q_0^2}{\zeta^2}C_\zeta.
\]
所以最后一个点的正号是两个负号相乘的结果.
全部新增展示式不另编号, 原文两条留数式及全部编号保持原样.
\subsection{两处谱无穷远的渐近展开}\label{ch:asymptotics}
谱曲面的两处无穷远分别为 $z=\infty$ 与 $z=0$. 以下在固定 $t$ 下推导修正 Jost 解及透射系数的展开. 基本条件特别蕴含
\begin{equation}\label{eq:highenergy-assumption}
q-q_0\in W^{2,1}(\mathbb R).
\end{equation}
此时 $q$ 有界, $q_x$ 连续且在两端趋于零, $q_x,q_{xx}\in L^1(\mathbb R)$.

\begin{proposition}[两处无穷远的渐近展开]\label{prop:highenergy}
在基本条件下, 记
\begin{equation}\label{eq:asymptotic-coefficients}
B(q)=\sg Q=\begin{pmatrix}0&q\\\conj{q}&0\end{pmatrix},\qquad
J_\pm(x)=\int_{\pm\infty}^{x}(q_0^2-|q(y)|^2)\,\dd y.
\end{equation}
各 Jost 解析列在相应闭域内满足
\begin{equation}\label{eq:mu-highenergy}
\mu_\pm(x,z)=I+\dfrac{A_\pm(x)}{z}+O(z^{-2}),\qquad
A_\pm=-\ii B(q)+\ii J_\pm\sg,\qquad z\to\infty,
\end{equation}
余项对实 $x$ 一致, 矩阵式按各自解析列理解. 令
\[
J_{\rm tot}=\int_{-\infty}^{+\infty}(q_0^2-|q|^2)\,\dd x,
\]
则透射系数满足
\begin{equation}\label{eq:transmission-highenergy}
\begin{aligned}
s_{22}(z)&=1+\dfrac{\ii J_{\rm tot}}{z}+O(z^{-2}),&&z\to\infty,\quad z\in\overline{D_+},\\
s_{11}(z)&=1-\dfrac{\ii J_{\rm tot}}{z}+O(z^{-2}),&&z\to\infty,\quad z\in\overline{D_-}.
\end{aligned}
\end{equation}
在另一处无穷远,
\begin{equation}\label{eq:smallmu}
\mu_\pm(x,z)=-\dfrac{\ii q_0}{z}\sigma_1
+\dfrac{1}{q_0}\begin{pmatrix}q(x)&-J_\pm(x)\\J_\pm(x)&\conj{q(x)}\end{pmatrix}
+O(z),\qquad z\to0,
\end{equation}
以及
\begin{equation}\label{eq:transmission-smallenergy}
\begin{aligned}
s_{22}(z)&=1+\dfrac{\ii J_{\rm tot}}{q_0^2}z+O(z^2),&&z\to0,\quad z\in\overline{D_+},\\
s_{11}(z)&=1-\dfrac{\ii J_{\rm tot}}{q_0^2}z+O(z^2),&&z\to0,\quad z\in\overline{D_-}.
\end{aligned}
\end{equation}
完整散射矩阵在实轴两端满足
\begin{equation}\label{eq:scattering-two-infinities}
\begin{aligned}
S(z)&=I-\dfrac{\ii J_{\rm tot}}{z}\sg+O(z^{-2}),&&z\to\infty,\\
S(z)&=I-\dfrac{\ii J_{\rm tot}}{q_0^2}z\sg+O(z^2),&&z\to0.
\end{aligned}
\end{equation}
\end{proposition}
\begin{proof}
先确定展开的系数. 将空间方程写为
\begin{equation}\label{eq:mu-highenergy-system}
\mu_x=-\dfrac{\ii z}{2}[\sg,\mu]+Q\mu
+\dfrac{\ii q_0^2}{2z}(\sg\mu+\mu\sg).
\end{equation}
以 $\mu=I+\dfrac{A}{z}+\dfrac{A_2}{z^2}+\cdots$ 为待证展开, 比较常数项与 $z^{-1}$ 项, 分别得到
\[
-\dfrac{\ii}{2}[\sg,A]+Q=0,\qquad
A_x=-\dfrac{\ii}{2}[\sg,A_2]+QA+\ii q_0^2\sg.
\]
第一式确定 $A$ 的非对角部分为 $-\ii\sg Q=-\ii B(q)$. 对第二式取对角部分, 使用 $QB=|q|^2\sg$, 得到
\[
(A_{\rm d})_x=\ii(q_0^2-|q|^2)\sg,\qquad
A_{\rm d}(\pm\infty)=0,
\]
故 $A_\pm=-\ii B+\ii J_\pm\sg$. 以下由积分方程证明此展开及其一致余项, 无须预先假定渐近级数可逐项求导.

作局部矩阵变换
\[
 T(x,z)=I-\dfrac{\ii B(q(x))}{z},\qquad \mu_\pm= T n_\pm.
\]
由 $B\sg=-\sg B$, $B^2=|q|^2I$,
\[
\det T=1+\dfrac{|q|^2}{z^2},\qquad
 T^{-1}=\dfrac{I+\dfrac{\ii B}{z}}{1+\dfrac{|q|^2}{z^2}}.
\]
当 $|z|$ 足够大时, $ T, T^{-1}$ 对所有实 $x$ 一致有界; 在归一化端 $ T\to Y$, 故 $n_\pm\to I$. 利用 $k+\lambda=z$ 和 $\lambda-k=\dfrac{q_0^2}{z}$, 有
\begin{align*}
U T+\ii\lambda T\sg- T_x
&=\ii(\lambda-k)\sg+Q-\dfrac{k+\lambda}{z}\sg B
-\dfrac{\ii}{z}QB+\dfrac{\ii}{z}B_x\\
&=\dfrac{\ii}{z}\bigl((q_0^2-|q|^2)\sg+B_x\bigr).
\end{align*}
因此
\begin{equation}\label{eq:n-volterra-system}
(n_\pm)_x=-\ii\lambda[\sg,n_\pm]+ Cn_\pm,\qquad
 C=\dfrac{1}{z} T^{-1} C_1,\qquad
 C_1=\ii\bigl((q_0^2-|q|^2)\sg+B_x\bigr).
\end{equation}
由 \eqref{eq:highenergy-assumption}, $q_0^2-|q|^2,q_x\in L^1$, 从而
\begin{equation}\label{eq:asymptotic-L1-control}
\| C(\cdot,z)\|_{L^1}\le\dfrac{C}{|z|},\qquad
\left\| C(\cdot,z)-\dfrac{ C_1}{z}\right\|_{L^1}
\le\dfrac{C}{|z|^2}.
\end{equation}
第 $j$ 列的 Volterra 方程为
\[
n_{\pm,j}(x,z)=e_j+\int_{\pm\infty}^{x}
e^{-\ii\lambda(x-y)(\sg-s_jI)} C(y,z)n_{\pm,j}(y,z)\,\dd y.
\]
其中 $s_1=1,s_2=-1$, 传播核在相应闭解析域及积分方向上范数不超过 1. 与定理 \ref{thm:jost} 相同的有序积分估计给出第 $l$ 次迭代项的界 $\dfrac{C^l}{l!|z|^l}$, 所以
\[
\|n_\pm-I\|_{L^\infty_x}\le e^{\frac{C}{|z|}}-1=O(z^{-1}),\qquad
\sum\limits_{l=2}^{\infty}\dfrac{C^l}{l!|z|^l}=O(z^{-2}).
\]
结合 \eqref{eq:asymptotic-L1-control}, 只需计算以 $\dfrac{ C_1}{z}$ 代替 $ C$ 的第一次迭代. 以负端第一列为例, 对角分量为
\[
\dfrac{1}{z}\int_{-\infty}^{x}( C_1)_{11}(y)\,\dd y
=\dfrac{\ii J_-(x)}{z},
\]
非对角分量为
\[
\dfrac{\ii}{z}\int_{-\infty}^{x}e^{2\ii\lambda(x-y)}\conj{q_y(y)}\,\dd y.
\]
由于 $q_x$ 在负端趋于零, 分部积分得
\[
\int_{-\infty}^{x}e^{2\ii\lambda(x-y)}\conj{q_y(y)}\,\dd y
=-\dfrac{\conj{q_x(x)}}{2\ii\lambda}
+\dfrac{1}{2\ii\lambda}\int_{-\infty}^{x}e^{2\ii\lambda(x-y)}\conj{q_{yy}(y)}\,\dd y.
\]
在 $\overline{D_+}$ 中, $y\le x$ 时指数模不超过 1, 且当 $z\to\infty$ 时 $|\lambda|\ge c|z|$. 因此该非对角分量绝对值至多为
\[
\dfrac{\|q_x\|_{L^\infty}+\|q_{xx}\|_{L^1}}{2|z\lambda|}=O(z^{-2}).
\]
正端和第二列只改变积分方向或指数符号, 具有相同估计. 于是
\[
n_\pm=I+\dfrac{\ii J_\pm\sg}{z}+O(z^{-2}),
\]
与 $ T$ 相乘即证明 \eqref{eq:mu-highenergy}. 此外, 从 \eqref{eq:n-volterra-system} 及 $n_\pm$ 非对角部分为 $O(z^{-2})$, 可得 $(n_\pm)_x=O(z^{-1})$, 因而 $(\mu_\pm)_x=O(z^{-1})$, 仍按解析列理解.

将展开代入 Wronskian 表达式, 有
\[
\det(\mu_{-,1},\mu_{+,2})
=1+\dfrac{\ii(J_--J_+)}{z}+O(z^{-2}),\qquad
J_--J_+=J_{\rm tot},\qquad \gamma=1+O(z^{-2}).
\]
这给出 $s_{22}$ 的展开. 同理,
\[
\det(\mu_{+,1},\mu_{-,2})
=1-\dfrac{\ii J_{\rm tot}}{z}+O(z^{-2}),
\]
从而得到 \eqref{eq:transmission-highenergy}. 在实轴上 $\lambda x$ 为实数, 由散射关系
\[
\begin{aligned}
S&=e^{\ii\lambda x\sg}\mu_-^{-1}\mu_+e^{-\ii\lambda x\sg},\\
\mu_-^{-1}\mu_+&=I+\dfrac{A_+-A_-}{z}+O(z^{-2})
=I-\dfrac{\ii J_{\rm tot}}{z}\sg+O(z^{-2}),
\end{aligned}
\]
得到 \eqref{eq:scattering-two-infinities} 的第一式.

最后, 令 $w=h(z)=-\dfrac{q_0^2}{z}$. 第 \ref{ch:symmetry} 节的换叶关系给出
\[
\mu_\pm(z)=-\dfrac{\ii q_0}{z}\mu_\pm(w)\sigma_1
=-\dfrac{\ii q_0}{z}\sigma_1+\dfrac{\ii}{q_0}A_\pm\sigma_1+O(z).
\]
计算 $\ii A_\pm\sigma_1=B\sigma_1-J_\pm\sg\sigma_1$ 即得 \eqref{eq:smallmu}. 因 $h$ 交换 $D_\pm$, 透射系数满足 $s_{22}(z)=s_{11}(w)$ 和 $s_{11}(z)=s_{22}(w)$; 将 $w^{-1}=-\dfrac{z}{q_0^2}$ 代入 \eqref{eq:transmission-highenergy}, 得到 \eqref{eq:transmission-smallenergy}. 在实轴上再用 $S(z)=\sigma_1S(w)\sigma_1$, 即得 \eqref{eq:scattering-two-infinities} 的第二式.
\end{proof}

两处无穷远控制 RHP 的归一化及光场重构. 分支点 $z=\pm\ii q_0$ 的连续延拓由引理 \ref{prop:jost-branch} 给出; 轮廓交点 $z=\pm q_0$ 是不同谱弧的交汇处, 共振原像 $z_a^\pm$ 则来自传播 Lax 矩阵的有理项.


\subsubsection*{两处无穷远展开的全部矩阵元与余项计算}
\paragraph{式 \eqref{eq:highenergy-assumption} 的具体用途.}
$q-q_0\in W^{2,1}$ 蕴含 $q-q_0,q_x,q_{xx}\in L^1$.
一维绝对连续代表满足
\[
q_x(x)=-\int_x^{+\infty}q_{yy}(y)\,\dd y,
\qquad
|q_x(x)|\le\|q_{xx}\|_{L^1}.
\]
同样从负端积分也可得到有界性与两端极限为零.
$q-q_0$ 亦有界, 所以 $q$ 有界.
对原文中对角积分的被积函数, 可以精确分解为
\begin{align*}
q_0^2-|q|^2
&=-(q-q_0)\conj{q}-q_0(\conj{q}-q_0),\\
|q_0^2-|q|^2|
&\le(|q|+q_0)|q-q_0|.
\end{align*}
因此该函数属于 $L^1$,
\[
|J_\pm(x)|\le(\|q\|_{L^\infty}+q_0)\|q-q_0\|_{L^1}
\]
对所有实 $x$ 一致成立.
$J_+$ 的正端定向要特别留意:
\[
J_-(x)=\int_{-\infty}^{x}(q_0^2-|q|^2)\,\dd y,
\qquad
J_+(x)=-\int_x^{+\infty}(q_0^2-|q|^2)\,\dd y.
\]
两者的空间导数相同,
\[
J_-'(x)=J_+'(x)=q_0^2-|q(x)|^2.
\]
但两者分别在各自归一化端为零.

\paragraph{式 \eqref{eq:mu-highenergy-system} 由原空间方程逐项得到.}
直接把 $k,\lambda$ 的有理表达式代入
$\mu_x=(-\ii k\sg+Q)\mu+\ii\lambda\mu\sg$,
\begin{align*}
\mu_x
&=-\dfrac{\ii z}{2}\sg\mu
+\dfrac{\ii q_0^2}{2z}\sg\mu
+Q\mu
+\dfrac{\ii z}{2}\mu\sg
+\dfrac{\ii q_0^2}{2z}\mu\sg\\
&=-\dfrac{\ii z}{2}(\sg\mu-\mu\sg)
+Q\mu+\dfrac{\ii q_0^2}{2z}(\sg\mu+\mu\sg).
\end{align*}
对易子与反对易部分逐元为
\[
[\sg,\mu]=\begin{pmatrix}0&2\mu_{12}\\-2\mu_{21}&0\end{pmatrix},
\qquad
\sg\mu+\mu\sg=\begin{pmatrix}2\mu_{11}&0\\0&-2\mu_{22}\end{pmatrix}.
\]
而势矩阵左乘的结果是
\[
Q\mu=\begin{pmatrix}
q\mu_{21}&q\mu_{22}\\
-\conj{q}\mu_{11}&-\conj{q}\mu_{12}
\end{pmatrix}.
\]
于是完整的四个标量微分方程为
\begin{align*}
(\mu_{11})_x&=q\mu_{21}+\dfrac{\ii q_0^2}{z}\mu_{11},\\
(\mu_{12})_x&=-\ii z\mu_{12}+q\mu_{22},\\
(\mu_{21})_x&=\ii z\mu_{21}-\conj{q}\mu_{11},\\
(\mu_{22})_x&=-\conj{q}\mu_{12}-\dfrac{\ii q_0^2}{z}\mu_{22}.
\end{align*}
从这四式可以逐个核对下面每一个系数的符号.

\paragraph{式 \eqref{eq:mu-highenergy}: 常数阶首先确定两个非对角元.}
按原文的待证形式把 $\mu=I+\dfrac{A}{z}+\dfrac{A_2}{z^2}+\cdots$ 代入.
左侧没有常数项, 右侧常数阶为
\[
-\dfrac{\ii}{2}[\sg,A]+Q
=\begin{pmatrix}
0&-\ii A_{12}+q\\
\ii A_{21}-\conj{q}&0
\end{pmatrix}=0.
\]
两个非平凡的标量等式解为
\[
-\ii A_{12}+q=0\ \Longrightarrow\ A_{12}=-\ii q,
\qquad
\ii A_{21}-\conj{q}=0\ \Longrightarrow\ A_{21}=-\ii\conj{q}.
\]
这个常数阶方程没有确定 $A_{11},A_{22}$,
因为与 $\sg$ 的对易子对角元恒为零.
不能把这两个尚未确定的对角系数直接设为零.

$z^{-1}$ 阶的矩阵方程中,
\[
-\dfrac{\ii}{2}[\sg,A_2]
=\begin{pmatrix}0&-\ii(A_2)_{12}\\
\ii(A_2)_{21}&0\end{pmatrix},
\qquad
QA=\begin{pmatrix}
qA_{21}&qA_{22}\\-\conj{q}A_{11}&-\conj{q}A_{12}
\end{pmatrix}.
\]
逐个比较四个矩阵元,
\begin{align*}
(A_{11})_x&=qA_{21}+\ii q_0^2,\\
(A_{12})_x&=-\ii(A_2)_{12}+qA_{22},\\
(A_{21})_x&=\ii(A_2)_{21}-\conj{q}A_{11},\\
(A_{22})_x&=-\conj{q}A_{12}-\ii q_0^2.
\end{align*}
将已经算出的 $A_{12},A_{21}$ 代入两条对角方程,
\begin{align*}
(A_{11})_x
&=q(-\ii\conj{q})+\ii q_0^2
=\ii(q_0^2-|q|^2),\\
(A_{22})_x
&=-\conj{q}(-\ii q)-\ii q_0^2
=-\ii(q_0^2-|q|^2).
\end{align*}
在归一化端 $\mu_\pm\to Y=I-\dfrac{\ii q_0}{z}\sigma_1$,
所以对角一阶系数的端态为零.
从相应端积分,
\begin{align*}
(A_\pm)_{11}(x)
&=\ii\int_{\pm\infty}^{x}(q_0^2-|q(y)|^2)\,\dd y
=\ii J_\pm(x),\\
(A_\pm)_{22}(x)
&=-\ii\int_{\pm\infty}^{x}(q_0^2-|q(y)|^2)\,\dd y
=-\ii J_\pm(x).
\end{align*}
因此四个系数完整合起来为
\[
A_\pm=\begin{pmatrix}
\ii J_\pm&-\ii q\\-\ii\conj{q}&-\ii J_\pm
\end{pmatrix}
=-\ii B(q)+\ii J_\pm\sg.
\]
这是系数的形式计算; 原文随后用积分方程证明余项,
所以不需要先假定整个渐近级数可以无限次逐项求导.

\paragraph{继续核对下一阶的两个非对角系数.}
虽然原结论只需一阶系数, 上面四个矩阵元还显示形式上的
$(A_2)_{12},(A_2)_{21}$ 应满足什么关系.
用 $A_{12}=-\ii q,A_{22}=-\ii J_\pm$,
\[
-\ii q_x=-\ii(A_2)_{12}-\ii qJ_\pm,
\qquad (A_2)_{12}=q_x-qJ_\pm.
\]
用 $A_{21}=-\ii\conj{q},A_{11}=\ii J_\pm$,
\[
-\ii\conj{q_x}=\ii(A_2)_{21}-\ii\conj{q}J_\pm,
\qquad (A_2)_{21}=-\conj{q_x}+\conj{q}J_\pm.
\]
这只是对待证级数的代数核对, 并不把原文的
$O(z^{-2})$ 余项擅自改成已证明的完整二阶展开.
原论证不需要对 $A_2$ 的对角部分再作积分.

\paragraph{原文局部矩阵 $ T$ 的逆逐元验证.}
这里仅使用原文已经引入的 $ T,n_\pm, C, C_1$,
不新增任何为简算而设的变换或辅助未知量.
\[
B=\begin{pmatrix}0&q\\\conj{q}&0\end{pmatrix},
\qquad
B^2=\begin{pmatrix}|q|^2&0\\0&|q|^2\end{pmatrix}=|q|^2I.
\]
于是
\[
 T=\begin{pmatrix}
1&-\dfrac{\ii q}{z}\\-\dfrac{\ii\conj{q}}{z}&1
\end{pmatrix},
\qquad
\det T
=1-\left(-\dfrac{\ii q}{z}\right)
\left(-\dfrac{\ii\conj{q}}{z}\right)
=1+\dfrac{|q|^2}{z^2}.
\]
将所写的逆矩阵与原矩阵相乘,
\begin{align*}
\left(I-\dfrac{\ii B}{z}\right)
\left(I+\dfrac{\ii B}{z}\right)
&=I+\dfrac{B^2}{z^2}
=\left(1+\dfrac{|q|^2}{z^2}\right)I.
\end{align*}
这说明逆矩阵公式中分母的符号为正.
只要 $|z|^2\ge2\|q\|_{L^\infty}^2$, 就有
\[
\left|1+\dfrac{|q(x)|^2}{z^2}\right|
\ge1-\dfrac{|q(x)|^2}{|z|^2}\ge\dfrac{1}{2}
\]
对所有 $x$ 一致成立, 因而逆矩阵一致有界.
在归一化端 $q\to q_0$, $ T\to Y$,
所以 $n_\pm= T^{-1}\mu_\pm\to I$.

\paragraph{式 \eqref{eq:n-volterra-system} 的剩余矩阵逐个矩阵元相消.}
先把 $U T$ 直接乘出来:
\[
U T=\begin{pmatrix}
-\ii k-\dfrac{\ii|q|^2}{z}&q-\dfrac{kq}{z}\\
-\conj{q}+\dfrac{k\conj{q}}{z}&\ii k+\dfrac{\ii|q|^2}{z}
\end{pmatrix}.
\]
另外两个矩阵分别为
\[
\ii\lambda T\sg
=\begin{pmatrix}
\ii\lambda&-\dfrac{\lambda q}{z}\\
\dfrac{\lambda\conj{q}}{z}&-\ii\lambda
\end{pmatrix},
\qquad
- T_x
=\begin{pmatrix}0&\dfrac{\ii q_x}{z}\\
\dfrac{\ii\conj{q_x}}{z}&0\end{pmatrix}.
\]
逐元相加, 左上元是
\[
-\ii k-\dfrac{\ii|q|^2}{z}+\ii\lambda
=\ii(\lambda-k)-\dfrac{\ii|q|^2}{z}
=\dfrac{\ii(q_0^2-|q|^2)}{z}.
\]
右下元是它的相反数.
右上元为
\[
q-\dfrac{(k+\lambda)q}{z}+\dfrac{\ii q_x}{z}
=q-q+\dfrac{\ii q_x}{z}
=\dfrac{\ii q_x}{z}.
\]
左下元为
\[
-\conj{q}+\dfrac{(k+\lambda)\conj{q}}{z}
+\dfrac{\ii\conj{q_x}}{z}
=-\conj{q}+\conj{q}+\dfrac{\ii\conj{q_x}}{z}
=\dfrac{\ii\conj{q_x}}{z}.
\]
因此剩余矩阵恰好为
\[
U T+\ii\lambda T\sg- T_x
=\dfrac{ C_1}{z},
\qquad
 C_1=\begin{pmatrix}
\ii(q_0^2-|q|^2)&\ii q_x\\
\ii\conj{q_x}&-\ii(q_0^2-|q|^2)
\end{pmatrix}.
\]

\paragraph{为什么 $n$ 方程中还有一个对易子.}
将 $\mu= Tn$ 代入原空间方程,
\[
 T_xn+ Tn_x
=U Tn+\ii\lambda Tn\sg.
\]
不能把最后一项直接改成 $\ii\lambda T\sg n$.
应加减同一项, 按原顺序写成
\begin{align*}
 Tn_x
&=(U T+\ii\lambda T\sg- T_x)n
+\ii\lambda T(n\sg-\sg n)\\
&=\dfrac{ C_1}{z}n
-\ii\lambda T[\sg,n].
\end{align*}
左乘 $ T^{-1}$ 就得到
$n_x=-\ii\lambda[\sg,n]+ Cn$.
整个推导中没有把非交换的矩阵乘法次序互换.

\paragraph{式 \eqref{eq:asymptotic-L1-control} 的两条估计.}
其中平方模差的可积性可逐项估计为
\begin{align*}
\left|q_0^2-|q|^2\right|
&=(q_0+|q|)\left|q_0-|q|\right|\\
&\le(q_0+\|q\|_{L^\infty})|q-q_0|,\\
\int_{-\infty}^{+\infty}\left|q_0^2-|q(x)|^2\right|\,\dd x
&\le(q_0+\|q\|_{L^\infty})\|q-q_0\|_{L^1}<\infty.
\end{align*}
第一行是平方差因式分解, 第二行使用反三角不等式.
由被积函数的可积性,
\[
\| C_1\|_{L^1}
\le C\left(\|q_0^2-|q|^2\|_{L^1}+\|q_x\|_{L^1}\right)<\infty.
\]
这里用于不等式上界的 $C>0$ 是与大谱参数无关的标量估计常数. 它与左侧的矩阵函数 $C(x,z)$ 及其系数 $C_1(x)$ 由自变量和维数区分. 第7节的 $T(x,z)$ 同样是二阶矩阵, 与第8节的标量指数 $T=e^{2s}$ 分别使用各自的定义.
原文的逆矩阵可改写成精确差值
\[
 T^{-1}-I
=\dfrac{\dfrac{\ii B}{z}-\dfrac{|q|^2}{z^2}I}
{1+\dfrac{|q|^2}{z^2}}.
\]
所以在前面的大 $|z|$ 范围,
\[
\| T^{-1}-I\|_{L^\infty_x}\le\dfrac{C}{|z|}.
\]
再按矩阵乘法范数估计,
\begin{align*}
\| C\|_{L^1}
&=\left\|\dfrac{1}{z} T^{-1} C_1\right\|_{L^1}
\le\dfrac{\| T^{-1}\|_{L^\infty}\| C_1\|_{L^1}}{|z|}
\le\dfrac{C}{|z|},\\
\left\| C-\dfrac{ C_1}{z}\right\|_{L^1}
&=\left\|\dfrac{1}{z}( T^{-1}-I) C_1\right\|_{L^1}\\
&\le\dfrac{\| T^{-1}-I\|_{L^\infty}\| C_1\|_{L^1}}{|z|}
\le\dfrac{C}{|z|^2}.
\end{align*}
因此不是把 $ T^{-1}$ 形式上替换为 $I$,
而是有一个可积的 $O(z^{-2})$ 误差.

\paragraph{四个 $n$ 列的传播核及第一次迭代.}
第一列对应
\[
e^{-\ii\lambda(x-y)(\sg-I)}
=\begin{pmatrix}1&0\\0&e^{2\ii\lambda(x-y)}\end{pmatrix},
\]
第二列对应
\[
e^{-\ii\lambda(x-y)(\sg+I)}
=\begin{pmatrix}e^{-2\ii\lambda(x-y)}&0\\0&1\end{pmatrix}.
\]
四种列与允许方向与第3节相同:
负端第一列和正端第二列用 $\operatorname{Im}\lambda\ge0$,
正端第一列和负端第二列用 $\operatorname{Im}\lambda\le0$.
故每个传播对角元的模不超过 1.
把 $\dfrac{ C_1}{z}$ 作用在自由项 $e_1,e_2$ 上,
第一次迭代分别为
\begin{align*}
n_{\pm,1}(x,z)
&=e_1+\dfrac{\ii}{z}
\begin{pmatrix}
\displaystyle\int_{\pm\infty}^{x}(q_0^2-|q|^2)\,\dd y\\[4pt]
\displaystyle\int_{\pm\infty}^{x}
 e^{2\ii\lambda(x-y)}\conj{q_y(y)}\,\dd y
\end{pmatrix}+O(z^{-2}),\\
n_{\pm,2}(x,z)
&=e_2+\dfrac{\ii}{z}
\begin{pmatrix}
\displaystyle\int_{\pm\infty}^{x}
 e^{-2\ii\lambda(x-y)}q_y(y)\,\dd y\\[4pt]
-\displaystyle\int_{\pm\infty}^{x}(q_0^2-|q|^2)\,\dd y
\end{pmatrix}+O(z^{-2}).
\end{align*}
这里的 $O(z^{-2})$ 包含系数替换误差与第二次以后的迭代.
具体地, 第 $l$ 次有序积分的上界为
$\dfrac{C^l}{l!|z|^l}$.
当 $|z|\ge1$ 时,
\begin{align*}
\sum\limits_{l=2}^{\infty}\dfrac{C^l}{l!|z|^l}
&\le\dfrac{1}{|z|^2}\sum\limits_{l=2}^{\infty}\dfrac{C^l}{l!}
=\dfrac{e^C-1-C}{|z|^2}.
\end{align*}
因此从第一次迭代开始分离余项是有定量依据的.

\paragraph{分部积分给出的额外一个大谱参数因子.}
对第一列的非对角积分,
\[
\partial_y e^{2\ii\lambda(x-y)}
=-2\ii\lambda e^{2\ii\lambda(x-y)}.
\]
把指数写成它的 $y$ 导数乘系数, 再分部积分,
\begin{align*}
&\int_{\pm\infty}^{x}e^{2\ii\lambda(x-y)}\conj{q_y(y)}\,\dd y\\
&=-\dfrac{1}{2\ii\lambda}
\int_{\pm\infty}^{x}\partial_y\bigl(e^{2\ii\lambda(x-y)}\bigr)
\conj{q_y(y)}\,\dd y\\
&=-\dfrac{1}{2\ii\lambda}
\left[e^{2\ii\lambda(x-y)}\conj{q_y(y)}\right]_{\pm\infty}^{x}
+\dfrac{1}{2\ii\lambda}
\int_{\pm\infty}^{x}e^{2\ii\lambda(x-y)}\conj{q_{yy}(y)}\,\dd y\\
&=-\dfrac{\conj{q_x(x)}}{2\ii\lambda}
+\dfrac{1}{2\ii\lambda}
\int_{\pm\infty}^{x}e^{2\ii\lambda(x-y)}\conj{q_{yy}(y)}\,\dd y.
\end{align*}
无穷远的边界项为零, 因为指数在当前允许方向有界且 $q_x\to0$.
第二列的指数导数符号相反,
\[
\partial_y e^{-2\ii\lambda(x-y)}
=2\ii\lambda e^{-2\ii\lambda(x-y)},
\]
所以逐项计算为
\begin{align*}
&\int_{\pm\infty}^{x}e^{-2\ii\lambda(x-y)}q_y(y)\,\dd y\\
&=\dfrac{1}{2\ii\lambda}
\left[e^{-2\ii\lambda(x-y)}q_y(y)\right]_{\pm\infty}^{x}
-\dfrac{1}{2\ii\lambda}
\int_{\pm\infty}^{x}e^{-2\ii\lambda(x-y)}q_{yy}(y)\,\dd y\\
&=\dfrac{q_x(x)}{2\ii\lambda}
-\dfrac{1}{2\ii\lambda}
\int_{\pm\infty}^{x}e^{-2\ii\lambda(x-y)}q_{yy}(y)\,\dd y.
\end{align*}
正端积分的方向仍为从 $+\infty$ 到 $x$;
定向积分的分部积分公式本身不需要另加符号.
两类积分的绝对值都不超过
\[
\dfrac{\|q_x\|_{L^\infty}+\|q_{xx}\|_{L^1}}{2|\lambda|}.
\]
对大 $|z|$,
\[
|\lambda(z)|
=\dfrac{1}{2}\left|z+\dfrac{q_0^2}{z}\right|
\ge\dfrac{1}{2}\left(|z|-\dfrac{q_0^2}{|z|}\right)
\ge\dfrac{|z|}{4}
\]
只需 $|z|^2\ge2q_0^2$.
所以第一次迭代前面的 $\dfrac{1}{z}$ 再乘这里的 $\dfrac{1}{\lambda}$,
确实给出 $O(z^{-2})$.

\paragraph{乘回 $ T$ 时四个矩阵元分别是什么.}
由对角积分与刚才的非对角估计,
\[
n_\pm=\begin{pmatrix}
1+\dfrac{\ii J_\pm}{z}+O(z^{-2})&O(z^{-2})\\
O(z^{-2})&1-\dfrac{\ii J_\pm}{z}+O(z^{-2})
\end{pmatrix}.
\]
直接把 $\mu_\pm= Tn_\pm$ 逐元相乘,
\begin{align*}
(\mu_\pm)_{11}
&=(n_\pm)_{11}-\dfrac{\ii q}{z}(n_\pm)_{21}
=1+\dfrac{\ii J_\pm}{z}+O(z^{-2}),\\
(\mu_\pm)_{12}
&=(n_\pm)_{12}-\dfrac{\ii q}{z}(n_\pm)_{22}
=-\dfrac{\ii q}{z}+O(z^{-2}),\\
(\mu_\pm)_{21}
&=-\dfrac{\ii\conj{q}}{z}(n_\pm)_{11}+(n_\pm)_{21}
=-\dfrac{\ii\conj{q}}{z}+O(z^{-2}),\\
(\mu_\pm)_{22}
&=-\dfrac{\ii\conj{q}}{z}(n_\pm)_{12}+(n_\pm)_{22}
=1-\dfrac{\ii J_\pm}{z}+O(z^{-2}).
\end{align*}
这里 $q,J_\pm$ 有界, 所以余项乘上这些系数仍有同一个阶.
四个余项在各自解析列允许的闭域与所有实 $x$ 上一致.

空间导数的估计也可以从 $n$ 方程逐元核对:
\[
[\sg,n_\pm]=\begin{pmatrix}0&O(z^{-2})\\O(z^{-2})&0\end{pmatrix}.
\]
因此 $\lambda[\sg,n_\pm]=O(z^{-1})$,
$ Cn_\pm=O(z^{-1})$, 故 $(n_\pm)_x=O(z^{-1})$.
同时 $ T_x=O(z^{-1})$, 所以
\[
(\mu_\pm)_x= T_xn_\pm+ T(n_\pm)_x=O(z^{-1}).
\]
这一步不是直接对没有先验证的 $O(z^{-2})$ 余项求导.

\paragraph{式 \eqref{eq:transmission-highenergy}: 两个 Wronskian 都展开.}
由四个矩阵元,
\begin{align*}
\det(\mu_{-,1},\mu_{+,2})
&=(\mu_-)_{11}(\mu_+)_{22}-(\mu_+)_{12}(\mu_-)_{21}\\
&=\left(1+\dfrac{\ii J_-}{z}+O(z^{-2})\right)
\left(1-\dfrac{\ii J_+}{z}+O(z^{-2})\right)\\
&\quad-\left(-\dfrac{\ii q}{z}+O(z^{-2})\right)
\left(-\dfrac{\ii\conj{q}}{z}+O(z^{-2})\right)\\
&=1+\dfrac{\ii(J_--J_+)}{z}+O(z^{-2}).
\end{align*}
非对角乘积至少是二阶, 所以不会改变一阶系数.
两个定向积分的差为
\begin{align*}
J_--J_+
&=\int_{-\infty}^{x}(q_0^2-|q|^2)\,\dd y
+\int_x^{+\infty}(q_0^2-|q|^2)\,\dd y\\
&=\int_{-\infty}^{+\infty}(q_0^2-|q|^2)\,\dd y=J_{\rm tot}.
\end{align*}
而
\[
\dfrac{1}{\gamma(z)}
=\dfrac{1}{1+\dfrac{q_0^2}{z^2}}
=1-\dfrac{q_0^2}{z^2}+O(z^{-4})
\]
没有一阶项, 所以 $s_{22}=1+\dfrac{\ii J_{\rm tot}}{z}+O(z^{-2})$.
另一个 Wronskian逐项为
\begin{align*}
\det(\mu_{+,1},\mu_{-,2})
&=(\mu_+)_{11}(\mu_-)_{22}-(\mu_-)_{12}(\mu_+)_{21}\\
&=\left(1+\dfrac{\ii J_+}{z}+O(z^{-2})\right)
\left(1-\dfrac{\ii J_-}{z}+O(z^{-2})\right)\\
&\quad-\left(-\dfrac{\ii q}{z}+O(z^{-2})\right)
\left(-\dfrac{\ii\conj{q}}{z}+O(z^{-2})\right)\\
&=1+\dfrac{\ii(J_+-J_-)}{z}+O(z^{-2})\\
&=1-\dfrac{\ii J_{\rm tot}}{z}+O(z^{-2}).
\end{align*}
除以同一个 $\gamma$ 就得到 $s_{11}$ 的展开.

\paragraph{两个非对角散射系数为什么没有一阶项.}
在实轴上, 两个 Wronskian 的一阶项直接抵消:
\begin{align*}
\det(\mu_{-,1},\mu_{+,1})
&=(\mu_-)_{11}(\mu_+)_{21}-(\mu_+)_{11}(\mu_-)_{21}\\
&=\left(1+\dfrac{\ii J_-}{z}\right)
\left(-\dfrac{\ii\conj{q}}{z}\right)
-\left(1+\dfrac{\ii J_+}{z}\right)
\left(-\dfrac{\ii\conj{q}}{z}\right)+O(z^{-2})\\
&=-\dfrac{\ii\conj{q}}{z}+\dfrac{\ii\conj{q}}{z}+O(z^{-2})
=O(z^{-2}),\\
\det(\mu_{+,2},\mu_{-,2})
&=(\mu_+)_{12}(\mu_-)_{22}-(\mu_-)_{12}(\mu_+)_{22}\\
&=\left(-\dfrac{\ii q}{z}\right)
\left(1-\dfrac{\ii J_-}{z}\right)
-\left(-\dfrac{\ii q}{z}\right)
\left(1-\dfrac{\ii J_+}{z}\right)+O(z^{-2})\\
&=-\dfrac{\ii q}{z}+\dfrac{\ii q}{z}+O(z^{-2})=O(z^{-2}).
\end{align*}
乘回 $e^{\pm2\ii(-\lambda x)}$ 后, 在实轴上指数模为 1,
所以 $s_{12},s_{21}=O(z^{-2})$.
这说明完整散射矩阵的一阶系数是对角矩阵.
在非实谱参数上, 不能不加条件地把指数因子的模也当作 1;
这正是原文把完整矩阵的展开限制在实轴上的原因.

\paragraph{式 \eqref{eq:scattering-two-infinities} 的矩阵计算交叉核对.}
由 $\det\mu_-=\gamma$, 二阶逆矩阵公式给出
\[
\mu_-^{-1}
=\dfrac{1}{\gamma}
\begin{pmatrix}(\mu_-)_{22}&-(\mu_-)_{12}\\
-(\mu_-)_{21}&(\mu_-)_{11}\end{pmatrix}
=I-\dfrac{A_-}{z}+O(z^{-2}).
\]
因而
\begin{align*}
\mu_-^{-1}\mu_+
&=\left(I-\dfrac{A_-}{z}+O(z^{-2})\right)
\left(I+\dfrac{A_+}{z}+O(z^{-2})\right)\\
&=I+\dfrac{A_+-A_-}{z}+O(z^{-2}),\\
A_+-A_-
&=-\ii B+\ii J_+\sg+\ii B-\ii J_-\sg\\
&=\ii(J_+-J_-)\sg=-\ii J_{\rm tot}\sg.
\end{align*}
对角系数与 $e^{-\ii\lambda x\sg}$ 交换,
所以左右相位共轭不改变它.
实轴上相位与逆相位有界, 因而也保留余项阶数,
最终得到原文 $S$ 在大 $z$ 的矩阵展开.

\paragraph{式 \eqref{eq:smallmu}: 用原换叶关系逐项乘出常数矩阵.}
这里的 $w=h(z)$ 是原文已经使用的谱点记号.
由于 $h(h(z))=z$, 在 $w$ 处应用换叶关系可得
\[
\mu_\pm(z)=\mu_\pm(w)C(w),
\qquad
C(w)=\dfrac{\ii w}{q_0}\sigma_1
=-\dfrac{\ii q_0}{z}\sigma_1.
\]
因此
\begin{align*}
\mu_\pm(z)
&=-\dfrac{\ii q_0}{z}
\left(I+\dfrac{A_\pm}{w}+O(w^{-2})\right)\sigma_1\\
&=-\dfrac{\ii q_0}{z}\sigma_1
+\dfrac{\ii}{q_0}A_\pm\sigma_1+O(z).
\end{align*}
这里使用 $w^{-1}=-\dfrac{z}{q_0^2}$;
大 $w$ 的 $O(w^{-2})$ 乘上 $\dfrac{1}{z}$ 后是 $O(z)$.
接着不省略常数矩阵乘法:
\[
B\sigma_1=\begin{pmatrix}q&0\\0&\conj{q}\end{pmatrix},
\qquad
\sg\sigma_1=\begin{pmatrix}0&1\\-1&0\end{pmatrix},
\]
\begin{align*}
\ii A_\pm\sigma_1
&=\ii(-\ii B+\ii J_\pm\sg)\sigma_1\\
&=B\sigma_1-J_\pm\sg\sigma_1\\
&=\begin{pmatrix}q&-J_\pm\\J_\pm&\conj{q}\end{pmatrix}.
\end{align*}
所以小 $z$ 的四个矩阵元分别为
\begin{align*}
(\mu_\pm)_{11}&=\dfrac{q}{q_0}+O(z),&
(\mu_\pm)_{12}&=-\dfrac{\ii q_0}{z}-\dfrac{J_\pm}{q_0}+O(z),\\
(\mu_\pm)_{21}&=-\dfrac{\ii q_0}{z}+\dfrac{J_\pm}{q_0}+O(z),&
(\mu_\pm)_{22}&=\dfrac{\conj{q}}{q_0}+O(z).
\end{align*}
非对角常数项的符号相反, 正是上面 $\sg\sigma_1$ 的两项符号不同.
在背景 $q=q_0,J_\pm=0$ 上, 这个式子恢复
$Y=I-\dfrac{\ii q_0}{z}\sigma_1$, 可直接复核常数矩阵.

\paragraph{式 \eqref{eq:transmission-smallenergy} 的每个符号.}
$h$ 交换两个解析域, 因而在 $z\in D_+$ 时 $w\in D_-$.
使用相应域内正确的透射展开,
\begin{align*}
s_{22}(z)
&=s_{11}(w)
=1-\dfrac{\ii J_{\rm tot}}{w}+O(w^{-2})\\
&=1-\ii J_{\rm tot}\left(-\dfrac{z}{q_0^2}\right)+O(z^2)
=1+\dfrac{\ii J_{\rm tot}}{q_0^2}z+O(z^2).
\end{align*}
在 $z\in D_-$ 时 $w\in D_+$, 同样有
\begin{align*}
s_{11}(z)
&=s_{22}(w)
=1+\dfrac{\ii J_{\rm tot}}{w}+O(w^{-2})\\
&=1+\ii J_{\rm tot}\left(-\dfrac{z}{q_0^2}\right)+O(z^2)
=1-\dfrac{\ii J_{\rm tot}}{q_0^2}z+O(z^2).
\end{align*}
小 $z$ 展开的正负号不仅由 $w^{-1}$ 的负号决定,
还使用了换叶后必须交换 $s_{11}$ 与 $s_{22}$.

完整散射矩阵在实轴上满足
\begin{align*}
S(z)
&=\sigma_1S(w)\sigma_1\\
&=\sigma_1\left(I-\dfrac{\ii J_{\rm tot}}{w}\sg
+O(w^{-2})\right)\sigma_1\\
&=I+\dfrac{\ii J_{\rm tot}}{w}\sg+O(w^{-2})\\
&=I-\dfrac{\ii J_{\rm tot}}{q_0^2}z\sg+O(z^2).
\end{align*}
第三行使用 $\sigma_1\sg\sigma_1=-\sg$,
第四行再使用 $w^{-1}=-\dfrac{z}{q_0^2}$.
两个非对角元也由换叶关系保持为 $O(z^2)$.
所有小 $z$ 公式都是从已验证的大 $z$ 展开与原约化关系推出,
没有在 $z=0$ 处对含奇系数的微分方程作不受控的形式除法.

\subsection{特殊谱点与后续 RHP 条件}
以上结论在基本条件下同时给出普通连续谱上的连续性、一阶矩支撑的分支点极限,
以及两处谱无穷远的展开. 三类有限特殊点的作用应分别说明.
在分支点 $\pm\ii q_0$, $Y$ 与 $\mu_\pm$ 的行列式退化;
在交点 $\pm q_0$, 空间方程与 $Y$ 正则, 但 RHP 的不同谱弧需要相容;
在共振原像 $z_a^\pm$, 空间 Jost 解连续可逆, 时间矩阵则具有固定极点.

Jost 列的连续延拓不自动给出除以透射系数后的有界性.
因此一般 RHP 中对分支点和交点的扇区边界值要求, 将在第 \ref{ch:rhp} 节单独列出.
无反射有理矩阵的分支点解析性与秩, 则由第 \ref{ch:soliton} 节的显式构造直接验证.
